\documentclass[11pt]{article}
\usepackage[margin=1in]{geometry}
\usepackage{amsmath,amssymb,amsthm,mathtools}
\usepackage[hidelinks,pdfusetitle]{hyperref}
\usepackage{microtype}
\usepackage{enumitem}
\usepackage{array}
\usepackage{booktabs}
\newcommand{\DOI}[1]{\mbox{\href{https://doi.org/#1}{\nolinkurl{#1}}}}
\newtheorem{theorem}{Theorem}[section]
\newtheorem{lemma}[theorem]{Lemma}
\newtheorem{corollary}[theorem]{Corollary}
\newtheorem{proposition}[theorem]{Proposition}
\newtheorem{observation}[theorem]{Observation}
\newtheorem{conjecture}[theorem]{Conjecture}
\newtheorem{obstruction}[theorem]{Route Obstruction}
\theoremstyle{definition}
\newtheorem{definition}[theorem]{Definition}
\theoremstyle{remark}
\newtheorem{remark}[theorem]{Remark}
\DeclareMathOperator{\Tr}{Tr}
\DeclareMathOperator{\spec}{spec}
\DeclareMathOperator{\dist}{dist}
\DeclareMathOperator{\diag}{diag}
\DeclareMathOperator{\ad}{ad}
\DeclareMathOperator{\Cov}{Cov}
\DeclareMathOperator{\Var}{Var}
\DeclareMathOperator{\sgn}{sgn}
\newcommand{\E}{\mathbb{E}}
\newcommand{\R}{\mathbb{R}}
\newcommand{\C}{\mathbb{C}}
\newcommand{\N}{\mathbb{N}}
\newcommand{\Z}{\mathbb{Z}}
\newcommand{\Tpair}{T_{\mathrm{pair}}}
\newcommand{\Nodd}{N_{\mathrm{odd}}}
\newcommand{\Dop}{\Delta}
\newcommand{\LF}{\mathrm{LF}}
\emergencystretch=1.5em
\title{Worldline-cluster, annealed-spectral, and telescoping obstructions to a phonon-frequency pairing ceiling in the Holstein model}
\author{Davis Rens\\SuperMatics Inc.\\\href{https://orcid.org/0000-0001-5839-5197}{ORCID 0000-0001-5839-5197}}
\date{}
\begin{document}
\maketitle

\begin{abstract}
Establishing an unconditional, volume- and chemical-potential-uniform pairing
ceiling $k_B \Tpair \le C(t,U)\Omega$ (up to possible logarithms) across all
couplings in the spinful Holstein model remains open, and this preprint does
not close it. Continuing three earlier obstruction papers
\cite{barriers1,barriers2,barriers3}, this fourth paper consolidates five
independently checked technical components that sharpen the boundary between
what holds rigorously and where current worldline, telescoping, and
annealed-field methods fail. First, in the worldline influence-functional
route, we prove the exact Parseval variance identity and disjoint-bump Gram
inequality for the Lang-Firsov phonon kernel together with an additive
multi-transition barrier lower bound (constant $3/22$, improving to $3/19$ for
half-width transitions), while identifying two structural obstructions at
$\beta t^2/U \gg 1$: signless $|\sigma|$-majorization fails on $16$ of $160$
four-hop neutral plaquette loops, and unexpanded-bath boundary subtraction
cannot cancel an $O(\beta t^2/U)$ extensive mismatch. Second, for synchronized
pair-hopping paths, we prove an exact telescoping obstruction capping bare
Lang-Firsov phonon suppression at a single $L$-independent endpoint factor and
Hubbard-Stratonovich determinant positivity for all $\mu \in \R$ with exact dimer
cancellation, and refute the literal strong-coupling bound of the preceding
framework via converged grand-canonical and canonical exact diagonalization.
Third, in the annealed static-plus-Gaussian framework, we prove the static
spectral package and a $(\beta\Omega)^2/96$ continuum first-variation bound,
establish the coherent-state pressure sandwich and quantum zero-mode variance
identity, rule out unrestricted low-temperature transfer of static level
counting via a dimer counterexample, and reject direct unweighted second-chaos
normalization into the static force-square. Fourth, we prove that an
unweighted annealed endpoint inequality gives only $\ell^1$ summability whereas
a spatially weighted exponential-moment closing estimate implies
$k_B \Tpair \le C(t,U)\Omega$; we prove the universal static Combes-Thomas base
bound $K_0 = (1+\pi/4)^2$, an exact $t=0$ barrier-to-well marginal suppression
that annihilates the pathwise staggered instanton, an unconditional $V=2$
operator-norm proof of the endpoint bound, and the exact $s=0$ derivative
bound $C_1^{\mathrm{phys}} = 1/24 + \zeta(3)/(8\pi^3)$. Fifth, we prove a
volume-, temperature-, frequency-, chemical-potential-, and background-uniform
Poghosyan-Ueltschi cluster expansion for the neutral-loop gas and its
perturbed exterior outside an $\Omega^{-1}$-collar, while delineating the open
interior collar screening $e^{+c U/\Omega}$ and four-hop virtual pair-return
channels $O(\beta t^4/(U^2\Omega))$.
\end{abstract}

\section{Introduction and program overview}
\label{sec:intro}
A central open problem in the theory of phonon-mediated superconductivity is
whether the pair-correlation decay temperature $\Tpair$ in the spinful
two-dimensional Holstein model~\cite{holstein1959} obeys a rigorous phonon-frequency ceiling
$k_B \Tpair \le C(t,U)\Omega$ (possibly up to logarithmic factors) uniformly in
the system volume and the chemical potential. Earlier released papers in this
program established an isotropic Migdal-Eliashberg ceiling
$k_B T < (\sqrt{377/250}/(2\pi))\sqrt{\rho}$~\cite{optimum} and its
finite-energy extension~\cite{finiteenergy}, an all-coupling sublinear ceiling
$k_B \Tpair \lesssim (U\Omega^2)^{1/3}\log^{1/3}$~\cite{linear4}, a
strong-coupling ceiling $k_B \Tpair \le 2^{27}t^2/U$ uniform in
$\Omega$~\cite{stage6} (building on~\cite{strong}), a cold-projected
complex-rotation bound~\cite{largeu}, and a Fisher-information framework for
phonon-mediated order~\cite{fisher}, alongside three earlier obstruction papers
\cite{barriers1,barriers2,barriers3} and the capstone synthesis~\cite{capstone}.
The most recent of these~\cite{barriers3} in particular isolated the discrepancy between
pathwise two-site Franck-Condon kernels and full-Gibbs cluster expansions,
leaving open whether a volume-uniform charged-skeleton influence bound or a
spectral-annealed second-chaos estimate could close the remaining gap in the
regime $\Omega \ll t^2/U$ when $\tau \sim \Omega$.

This fourth paper consolidates five independently checked technical packets that
map out what can be proved on both sides of that gap and where the leading
closure strategies break down. Every theorem, lemma, proposition, and corollary
stated below has a complete derivation, while every route obstruction is backed
by an explicit mathematical witness or counterexample. Across the five
sections, we establish:
\begin{enumerate}[leftmargin=1.5em]
\item \emph{Charged-skeleton influence bounds and route obstructions}
(Section~\ref{sec:charged-skeleton}): an exact Parseval identity and
disjoint-bump Gram inequality for the Lang-Firsov phonon action $\Phi_\beta[q_z]$
in the full unrestricted Fock space, an additive multi-transition barrier lower
bound with constant $3/22$ (improving to $3/19$ for half-width transitions and
robust to short opposite-charge blips), a comparison with the
$\tau \ge 6000\,t^2/U$ regime~\cite{strong}, and two structural route
obstructions at $\beta t^2/U \gg 1$ (including $16$ negative four-hop neutral
plaquette loops out of $160$).
\item \emph{Telescoping obstruction, determinant positivity, and exact
diagonalization refutation} (Section~\ref{sec:telescope}): exact Lang-Firsov
telescoping capping synchronized bare-frame suppression at a single
$L$-independent endpoint factor and collapsing to $1$ on closed physical
correlators, generalized Hubbard-Stratonovich determinant positivity on every
finite graph and real chemical potential $\mu$ with exact two-site cancellation,
and a numerical refutation of the literal strong-coupling physical-correlator
bound with $O(1)$ neutral remainder.
\item \emph{Annealed spectral estimates and the second-chaos obstruction}
(Section~\ref{sec:annealed-spectral}): the static annealed Fermi-kernel,
relative-entropy level-counting, Cooper-bubble, and Fisher-integration
package, a $(\beta\Omega)^2/96$ continuum first-variation bound, the full
coherent-state pressure sandwich, the exact quantum zero-mode variance
identity, a low-temperature dimer counterexample ruling out unrestricted
transfer of static level counting to the quantum measure, and a negative
verdict on direct unweighted second-chaos normalization into the static
force-square.
\item \emph{Annealed endpoint reduction, static Combes-Thomas base, and
instanton annihilation} (Section~\ref{sec:annealed-endpoint}): a conditional
reduction showing why an unweighted endpoint expectation yields only $\ell^1$
summability while a spatially weighted exponential-moment closing estimate
yields $k_B \Tpair \le C(t,U)\Omega$, the universal static Combes-Thomas base
bound $K_0 = (1+\pi/4)^2$ (plus an interval-certified $V=6$ Schur test), an
exact $t=0$ nine-component marginal and barrier-to-well density ratio
$R_0(B;\beta,U)$ that annihilates the staggered instanton, an unconditional
$V=2$ operator-norm proof of the endpoint bound, the exact $s=0$ derivative
constant $C_1^{\mathrm{phys}} = 1/24 + \zeta(3)/(8\pi^3)$, and the
$\Omega \to \infty$ attractive Hubbard anchor together with two remaining
closure barriers.
\item \emph{Volume-uniform neutral-dipole cluster expansion and collar
boundary} (Section~\ref{sec:neutral-dipole}): a volume-, temperature-,
frequency-, chemical-potential-, and background-uniform Poghosyan-Ueltschi
cluster expansion for the projected neutral-loop gas and its perturbed exterior
outside an $\Omega^{-1}$-collar around any charged skeleton, with explicit
constants across degrees $\Delta \in \{2,4,6\}$, while isolating the remaining
open $e^{+cU/\Omega}$ interior collar screening and four-hop virtual
pair-return channels $O(\beta t^4/(U^2\Omega))$.
\end{enumerate}

Throughout the paper we set $\hbar = k_B = 1$, write $\beta = \tau^{-1} = T^{-1}$,
$X = U/\Omega$, and $\alpha = \sqrt{X/2}$, and preserve the native parameter
conventions of each component explicitly at the start of its section so that
each statement matches its verified proof and replay script. The target global
bound $k_B \Tpair \le C(t,U)\Omega$ remains open.

\begin{definition}[Pair-correlation decay temperature; \cite{largeu,strong,barriers3}]
\label{def:tpair}
For fixed $t \ge 0$, $U > 0$, and $\Omega > 0$, let $\Tpair$ denote the infimum
of temperatures $\tau = k_B T > 0$ such that for every $\tau' > \tau$ there
exist constants $A_0 > 0$ and $\xi > 0$, uniform in the system size $L \ge 4$
and the chemical potential $\mu \in \R$, satisfying
\begin{equation}
\label{eq:tpair-def}
\bigl|\langle \Dop_x^\dagger \Dop_y \rangle\bigr| \le A_0 e^{-d(x,y)/\xi}
\qquad \forall x, y \in \Lambda,
\end{equation}
and $\sup_{L \ge 4} \max_x \sum_y |\langle \Dop_x^\dagger \Dop_y \rangle| < \infty$.
As in~\cite{barriers3}, $\Tpair$ is an upper threshold for uniform exponential
decay and row summability of the singlet pair correlator, not a critical
temperature of a phase transition.
\end{definition}

\section{Charged-skeleton influence bounds and route obstructions}
\label{sec:charged-skeleton}

Let $\Lambda = (\Z/L\Z)^2$ ($L \ge 4$, periodic torus) with edge set $E(\Lambda)$,
hopping $t \ge 0$, optical phonon frequency $\Omega > 0$, attraction $U > 0$,
and inverse temperature $\beta = \tau^{-1} > 0$, with $x = t/U$, $X = U/\Omega$,
and $\alpha = \sqrt{X/2}$. On the full unrestricted Fock space
$\mathcal{H}_\Lambda = \bigotimes_{z \in \Lambda} (\C^4 \otimes \ell^2(\Z_{\ge 0}))$
(with no phonon-number cutoff and no cold projector), the spinful Holstein
Hamiltonian in the pre-Lang-Firsov frame is
\begin{equation}
\label{eq:cs-Hpre}
H_\Lambda^{\mathrm{pre}}
= \Omega \sum_{z \in \Lambda} a_z^\dagger a_z
- \sqrt{\frac{U\Omega}{2}} \sum_{z \in \Lambda} (n_z - 1)(a_z^\dagger + a_z)
- t \sum_{\langle zw \rangle \in E(\Lambda), \sigma \in \{\uparrow,\downarrow\}}
  \bigl(c_{z\sigma}^\dagger c_{w\sigma} + c_{w\sigma}^\dagger c_{z\sigma}\bigr)
- \mu \sum_{z \in \Lambda} (n_z - 1),
\end{equation}
where $n_z = n_{z\uparrow} + n_{z\downarrow}$ and $\Dop_z = c_{z\downarrow} c_{z\uparrow}$.
Conjugating by the Lang-Firsov unitary $\mathcal{U} = \exp[-\alpha \sum_z (n_z - 1)(a_z^\dagger - a_z)]$~\cite{langfirsov}
yields $H_{\LF} = \mathcal{U} H_\Lambda^{\mathrm{pre}} \mathcal{U}^\dagger = H_0 + V - \frac{U}{2}|\Lambda|$,
where
\[
H_0 = \Omega \sum_{z \in \Lambda} a_z^\dagger a_z + \frac{U}{2}\Nodd - \mu \sum_{z \in \Lambda}(n_z - 1),
\qquad
\Nodd := \sum_{z \in \Lambda} \mathbf{1}_{\{n_z = 1\}},
\]
and $V = -t \sum_{\langle zw\rangle, \sigma} (c_{z\sigma}^\dagger c_{w\sigma} D_z(-\alpha) D_w(\alpha) + \mathrm{h.c.})$
with unitary displacements $D_z(\gamma) = e^{\gamma(a_z^\dagger - a_z)}$. Because
$\|V\|_{\mathrm{op}} \le 4|E(\Lambda)|t < \infty$ and $e^{-\beta H_0}$ is trace-class
with $Z_{\mathrm{ph}}(\beta) := (1 - e^{-\beta\Omega})^{-1}$, the Duhamel series
for $Z_\Lambda = \Tr_{\mathcal{H}_\Lambda}(e^{-\beta H_{\LF}})$ and
$Z_\Lambda^{xy} = \Tr_{\mathcal{H}_\Lambda}(\mathcal{U}\Dop_x^\dagger \Dop_y \mathcal{U}^\dagger e^{-\beta H_{\LF}})$
converges absolutely in trace norm for every finite $\Lambda$, $\beta > 0$,
and $\mu \in \R$. Integrating out the harmonic oscillators produces the exact
continuous-time worldline representation~\cite{barriers3}:
\[
Z_\Lambda = Z_{\mathrm{ph}}(\beta)^{|\Lambda|}
  \sum_{\omega \in \mathcal{W}_\Lambda} \sigma(\omega)\, w_{\mathrm{el}}(\omega)\, e^{-\Phi_\beta[\omega]},
\qquad
Z_\Lambda^{xy} = Z_{\mathrm{ph}}(\beta)^{|\Lambda|}
  \sum_{\omega \in \mathcal{W}_\Lambda^{xy}} \sigma(\omega)\, w_{\mathrm{el}}(\omega)\, e^{-\Phi_\beta[\omega]},
\]
where $q_z(u) := n_z(u) - 1 \in \{-1,0,+1\}$, $\sigma(\omega) \in \{\pm 1\}$ is
the fermionic sign, $T_{\mathrm{odd}}[q_z] = |\{u \in S^1_\beta : q_z(u) = 0\}|$,
$w_{\mathrm{el}}(\omega) = t^{h(\omega)} e^{-\frac{U}{2} T_{\mathrm{odd}}(\omega) + \mu \int_0^\beta \sum_z q_z(u)\,du}\prod du_j \ge 0$,
and $\Phi_\beta[\omega] = \sum_{z \in \Lambda} \Phi_\beta[q_z]$ with
\begin{equation}
\label{eq:cs-phi}
\Phi_\beta[q] = \frac{1}{2}\alpha^2 \sum_{j,k} c_j c_k D_\beta(u_j - u_k) \ge 0,
\qquad
D_\beta(u) = \frac{\cosh(\Omega(\beta/2 - |u|_\beta))}{\sinh(\beta\Omega/2)},
\end{equation}
for a piecewise constant periodic path $q: S^1_\beta \to \{-1,0,+1\}$ with jumps
$\{c_k\}$ at $\{u_k\}$ ($\sum_k c_k = 0$) and circle distance
$|u|_\beta := \min_{n \in \Z}|u - n\beta| \in [0,\beta/2]$.

\begin{lemma}[Parseval identity and disjoint-bump Gram inequality]
\label{lem:cs-gram}
Let $q: S^1_\beta \to \{-1,0,+1\}$ be a finite-jump periodic charge path on
$S^1_\beta = \R/\beta\Z$ with Fourier coefficients
$\hat{q}_m = \frac{1}{\beta}\int_0^\beta q(u) e^{-i\omega_m u}\,du$
($\omega_m = 2\pi m/\beta$). Then:
\begin{enumerate}[leftmargin=1.5em]
\item \emph{(Exact Parseval variance identity)}
\[
\Phi_\beta[q]
= \frac{U\beta}{2} \sum_{m \in \Z \setminus \{0\}}
  \frac{\omega_m^2}{\omega_m^2 + \Omega^2} |\hat{q}_m|^2
= \frac{U}{2} \|q - \bar{q}\|_{L^2(S^1_\beta)}^2
- \frac{U}{2} \langle q - \bar{q},\, K_\Omega (q - \bar{q})\rangle_{L^2(S^1_\beta)},
\]
where $\bar{q} = \frac{1}{\beta}\int_0^\beta q(u)\,du$,
$\widehat{(K_\Omega f)}_m = \frac{\Omega^2}{\omega_m^2 + \Omega^2}\hat{f}_m$
for $m \ne 0$, and $K_\Omega 1 = 0$.
\item \emph{(Exact weak derivative form)} For every $\beta$-periodic
$g \in H^1(S^1_\beta)$ with zero mean,
\begin{equation}
\label{eq:cs-parseval-h1}
\Phi_\beta[q]
\ge \frac{U}{2} \frac{\bigl|\int_0^\beta q(u) g'(u)\,du\bigr|^2}{\int_0^\beta \bigl(g'(u)^2 + \Omega^2 g(u)^2\bigr)\,du}.
\end{equation}
\item \emph{(Disjoint-bump Gram inequality)} If $g_1,\dots,g_m \in H^1(S^1_\beta)$
have mutually disjoint open supports modulo $\beta$, then the $m \times m$
Gram matrix $G_{ab} = \langle g_a', g_b'\rangle_{L^2} + \Omega^2 \langle g_a - \bar{g}_a, g_b - \bar{g}_b\rangle_{L^2}$
in the norm $\|f\|_{\Omega}^2 = \beta \sum_{n \ne 0}(\omega_n^2 + \Omega^2)|\hat{f}_n|^2$
satisfies $G \le \diag(d_1,\dots,d_m)$ in the positive-semidefinite order, where
$d_k = \int_0^\beta (g_k'(u)^2 + \Omega^2 g_k(u)^2)\,du$, and consequently
\[
\Phi_\beta[q]
\ge \frac{U}{2} \sum_{k=1}^m \frac{\bigl(\int_0^\beta q(u) g_k'(u)\,du\bigr)^2}{\int_0^\beta \bigl(g_k'(u)^2 + \Omega^2 g_k(u)^2\bigr)\,du}.
\]
\end{enumerate}
\end{lemma}

\begin{proof}
Writing $c_j$ as the distributional derivative $dq = \sum_j c_j \delta_{u_j} du$,
the Fourier coefficients of $dq$ are $i\omega_m \beta \hat{q}_m$, while
$\hat{D}_\beta(\omega_m) = \frac{2\Omega}{\omega_m^2 + \Omega^2}$ gives
$\Phi_\beta[q] = \frac{\alpha^2 \Omega}{\beta}\sum_{m \ne 0} \frac{\omega_m^2 \beta^2}{\omega_m^2 + \Omega^2}|\hat{q}_m|^2$,
which simplifies via $\alpha^2 \Omega = U/2$ to the Parseval identity in~(1).
For a real zero-mean test function $g \in H^1(S^1_\beta)$,
Cauchy-Schwarz in $\ell^2(\Z \setminus \{0\})$ with weights
$\frac{\omega_m^2}{\omega_m^2 + \Omega^2}$ and $\omega_m^2 + \Omega^2$ gives
$|\langle q, g'\rangle_{L^2}|^2 = \beta^2 |\sum_{m\ne 0} i\omega_m \hat{q}_m \overline{\hat{g}_m}|^2 \le (\frac{2}{U}\Phi_\beta[q])\cdot \int_0^\beta (g'^2 + \Omega^2 g^2)\,du$,
yielding~\eqref{eq:cs-parseval-h1}. For mutually disjointly supported bumps
$g_1,\dots,g_m$, the $L^2$ inner products $\langle g_a', g_b'\rangle_{L^2}$ and
$\langle g_a, g_b\rangle_{L^2}$ vanish for $a \ne b$, while subtracting the zero-mean
projection in the dual norm $\beta \sum_{n \ne 0}(\omega_n^2 + \Omega^2)|\hat{f}_n|^2$
only lowers $G$ by the rank-one positive-semidefinite matrix
$\beta\Omega^2 \bar{g}\bar{g}^T$ (with $\bar{g}_a = \frac{1}{\beta}\int_0^\beta g_a(u)\,du$),
giving $G \le \diag(d_1,\dots,d_m)$. Applying Bessel's inequality to
$v_k = \int_0^\beta q g_k'\,du$ against $G^{-1}$ on $\operatorname{ran}(G)$
yields~(3).
\end{proof}

Lemma~\ref{lem:cs-gram} converts disjoint spacetime transitions into an additive
lower bound on the combined phonon and atomic cost.

\begin{lemma}[Additive multi-transition barrier lower bound]
\label{lem:cs-barrier}
Let $\ell_0 := \min(1/\Omega, \beta/4)$ and let $z \in \Lambda$ have $m \ge 1$
disjoint transition windows $S_k = A_k^- \sqcup G_k \sqcup A_k^+ \subset S^1_\beta$
($k=1,\dots,m$), each of total length $|S_k| = 2\ell_0 + \delta_k$, with
$|A_k^-| = |A_k^+| = \ell_0$, $|G_k| = \delta_k \ge 0$, and signs
$\varepsilon_k \in \{\pm 1\}$ such that:
\begin{itemize}[leftmargin=1.5em]
\item the minority-charge blip measure on the flanks satisfies
$|\{u \in A_k^- : q_z(u) = +\varepsilon_k\}| \le \frac{1}{8}\ell_0$ and
$|\{u \in A_k^+ : q_z(u) = -\varepsilon_k\}| \le \frac{1}{8}\ell_0$;
\item either $\delta_k \le \ell_0$, or $q_z \equiv 0$ on $G_k$.
\end{itemize}
Then, uniformly over all neutral excursions elsewhere on $S^1_\beta$ and on
all other sites $\Lambda \setminus \{z\}$,
\begin{equation}
\label{eq:cs-barrier}
\Phi_\beta[q_z] + \frac{U}{2} T_{\mathrm{odd}}[q_z]
\ge \frac{3}{22}\, m\, U \ell_0
= \frac{3}{22}\, m\, X \min\!\Bigl(1, \frac{\beta\Omega}{4}\Bigr).
\end{equation}
If $\delta_k \le \frac{1}{2}\ell_0$ for every $k$, the constant $3/22$ improves
to $3/19$.
\end{lemma}

\begin{proof}
Partition $\{1,\dots,m\} = \mathcal{I}_1 \sqcup \mathcal{I}_2$, placing $k$ in
$\mathcal{I}_1$ if $T_{\mathrm{odd},k} := |\{u \in S_k : q_z(u) = 0\}| \ge \frac{1}{2}\ell_0$
and in $\mathcal{I}_2$ otherwise. For $k \in \mathcal{I}_1$, the atomic term on
$S_k$ alone contributes $\frac{U}{2} T_{\mathrm{odd},k} \ge \frac{1}{4}U\ell_0 > \frac{3}{22}U\ell_0$.
For $k \in \mathcal{I}_2$, necessarily $\delta_k \le \ell_0$ (otherwise $q_z \equiv 0$
on $G_k$ would put $k \in \mathcal{I}_1$). Choose the piecewise linear bump $g_k$
supported on $S_k$ that ramps $0 \to 1$ on $A_k^-$, equals $1$ on $G_k$, and
ramps $1 \to 0$ on $A_k^+$. On $A_k^-$, $q_z = -\varepsilon_k$ on a set of
measure at least $\ell_0 - T_{\mathrm{odd},k}^- - \frac{1}{8}\ell_0$ and
$q_z = +\varepsilon_k$ on at most $\frac{1}{8}\ell_0$, and symmetrically on
$A_k^+$, giving
\[
\Bigl|\int_{S_k} q_z g_k'\,du\Bigr|
\ge \frac{1}{\ell_0}\Bigl(2\ell_0 - T_{\mathrm{odd},k} - \frac{1}{2}\ell_0\Bigr)
= \frac{3}{2} - \frac{T_{\mathrm{odd},k}}{\ell_0} > 1.
\]
Meanwhile $\int_{S_k} ((g_k')^2 + \Omega^2 g_k^2)\,du \le \frac{2}{\ell_0} + \Omega^2 (\frac{2}{3}\ell_0 + \delta_k) \le \frac{11}{3\ell_0}$
(using $\Omega\ell_0 \le 1$ and $\delta_k \le \ell_0$; when $\delta_k \le \frac{1}{2}\ell_0$,
this tightens to $\frac{19}{6\ell_0}$). By Lemma~\ref{lem:cs-gram}(3) and
$T_{\mathrm{odd}}[q_z] \ge \sum_k T_{\mathrm{odd},k}$, setting
$s_k := T_{\mathrm{odd},k}/\ell_0 \in [0, 1/2)$ gives the per-window contribution
$\frac{U\ell_0}{2} [\frac{3}{11}(3/2 - s_k)^2 + s_k]$. Differentiating with respect to
$s_k$ shows that this quadratic is increasing on $[0, 1/2)$ (its derivative at
$s_k = 0$ is $1 - \frac{9}{11} = \frac{2}{11} > 0$), so its minimum on
$[0, 1/2]$ occurs at $s_k = 0$ with value $\frac{27}{88} U\ell_0 > \frac{3}{22} U\ell_0$
(and $\frac{27}{76} U\ell_0 > \frac{3}{19} U\ell_0$ when $\delta_k \le \ell_0/2$;
explicitly, optimizing over $s_k \in [0, 1/2]$ with the tighter constant gives
at least $\frac{3}{22}$ in general and $\frac{3}{19}$ for $\delta_k \le \ell_0/2$,
which also holds for $k \in \mathcal{I}_1$).
\end{proof}

\begin{remark}[Arithmetic parameter window and comparison with prior strong-coupling bounds]
\label{rem:cs-window}
Comparing the single-bridge barrier $\frac{3}{22} X \min(1, \beta\Omega/4)$ of
Lemma~\ref{lem:cs-barrier} against the worst-case neutral Ising energy change
per flipped bond of order $\beta t^2/U$ shows that the energetic balance
$\beta t^2/U \le \eta_0 X \min(1,\beta\Omega/4)$ holds arithmetically even in
the regime $\tau \sim \Omega \ll t^2/U$ whenever $(\beta\Omega)(t/U)^2 \ll 1$.
For $\beta t^2/U \le 1/6000$, the homogeneous full-Gibbs ceiling
$|\langle \Dop_x^\dagger \Dop_y\rangle_\beta| \le \frac{57}{500}(99/100)^{d(x,y)}$
is already established in~\cite{strong}. Extending a volume-uniform ceiling into
$\beta t^2/U \gg 1$, however, hits two structural obstructions, consistent with
the exact two-site grand-canonical value $\langle \Dop_0^\dagger \Dop_1 \rangle_\beta \approx 0.1376(t/U)^2$
at $U=20,\Omega=1,t=0.5,\tau=1.0$ exceeding the bare single-bridge kernel
$K_\beta^{\mathrm{bare}}(1) \approx 8.58 \times 10^{-38}$ by more than thirty
orders of magnitude without refuting a bound that leaves a neutral prefactor
unconstrained.
\end{remark}

\begin{obstruction}[Neutral four-hop fermion sign witness]
\label{obs:cs-sign}
On a $2 \times 2$ square plaquette starting from the checkerboard state
$\eta = (1,0,1,0)$ (pairs on $(0,0)$ and $(1,1)$, empty sites $(1,0)$ and $(0,1)$),
there exist $16$ distinct four-hop neutral plaquette exchange loops out of $160$
total closed four-hop neutral loops that carry fermionic sign $\sigma(\omega) = -1$
and weight $\Theta(\beta t^4/U^3)$. Consequently, replacing $\sigma(\omega)$ by
$|\sigma(\omega)| = 1$ in the neutral background introduces an uncompensated
plaquette discrepancy of order $\beta t^4/U^3$ per plaquette, obstructing
naive signless $|\sigma|$-majorization as $|\Lambda| \to \infty$ whenever
$|\Lambda|\beta t^4/U^3 \gtrsim 1$. (By contrast, spacetime-disjoint neutral
polymers $\gamma_1, \gamma_2$ satisfy $\sigma(\gamma_1 \cup \gamma_2) = \sigma(\gamma_1)\sigma(\gamma_2)$.)
\end{obstruction}

\begin{obstruction}[Unexpanded neutral-loop boundary coupling at $\beta t^2/U \gg 1$]
\label{obs:cs-boundary}
If one instead retains the exact positive operator trace
$Z_{\mathrm{neut}}(\eta) = \Tr_{P_\eta \mathcal{H}}(e^{-\beta H_{\mathrm{neut}}(\eta)}) > 0$,
inserting a charged polymer with spatial support $S$ that flips $\eta_S \mapsto \eta_S'$
leaves the cluster coupled across the $2d|S|$ boundary edges $\partial S$ to an
unexpanded neutral bath in $\Lambda \setminus S$. Because the boundary hopping
has operator norm $O(|S|t)$ rather than $O(|S|t^2/U)$, avoiding a $\exp(C\beta t)$
loss at $\tau \ll (t/U)\Omega$ requires a volume-uniform bulk subtraction and
interface-decoration theorem separating the $N_{\mathrm{odd}}(\Lambda \setminus S) \ge 1$
sector on arbitrary excited configurations $\eta_{\Lambda\setminus S}$.
\end{obstruction}
\section{Telescoping obstruction, determinant positivity, and exact diagonalization refutation}
\label{sec:telescope}

Let $V$ be a finite nearest-neighbor graph with edge set $E$ and adjacency
matrix $A_V$, and set $D_\beta(s) = \cosh(\Omega(\beta/2 - |s|_\beta))/\sinh(\beta\Omega/2)$
with $|s|_\beta = \min(|s|, \beta - |s|)$. Under the Lang-Firsov unitary
$\mathcal{U} = \exp[-\alpha \sum_z (n_z - 1)(a_z^\dagger - a_z)]$,
$\widetilde{\Dop}_z := \mathcal{U}\Dop_z\mathcal{U}^\dagger = \Dop_z D_z(+2\alpha)$
and $\widetilde{\Dop}_z^\dagger = \Dop_z^\dagger D_z(-2\alpha)$, where
$D_z(r) = \exp[r(a_z^\dagger - a_z)]$. Following the two-site kernel convention
of~\cite{barriers3}, define
\[
F_-(u) := \exp\!\left[-\frac{U}{\Omega}\bigl(D_\beta(0) - D_\beta(u)\bigr)\right],
\qquad
F_-(u)^2 = \exp\!\left[-\frac{2U}{\Omega}\bigl(D_\beta(0) - D_\beta(u)\bigr)\right],
\]
so that in the single-site convention $F_-^{\mathrm{FG}}(u)$ used earlier,
$F_-^{\mathrm{FG}}(u)^4 = F_-(u)^2$.

\begin{theorem}[Exact synchronized pair-transfer telescoping]
\label{thm:tel-telescope}
Let $\gamma = (x_0, x_1, \dots, x_L)$ be any walk of length $L \ge 1$ on $V$,
with synchronized pair-transfer times $0 < s_1 < s_2 < \dots < s_L < \beta$
(at each $s_j$, both spins hop $x_{j-1} \to x_j$, contributing
$D_{x_{j-1}}(+2\alpha) D_{x_j}(-2\alpha)$).
\begin{enumerate}[leftmargin=1.5em]
\item \emph{(Open bare-LF trace at coincidence)} As $s_1, \dots, s_L \to s \in (0,\beta)$
without endpoint observables, $\prod_{j=1}^L D_{x_{j-1}}(+2\alpha) D_{x_j}(-2\alpha) = D_{x_0}(+2\alpha) D_{x_L}(-2\alpha)$.
For $x_0 \neq x_L$, the exact thermal phonon trace is
\[
W_{\mathrm{open},\,\mathrm{sync}}(\gamma; s,\dots,s)
= \exp\!\left[-2 X \coth\!\left(\frac{\beta\Omega}{2}\right)\right],
\]
strictly independent of $L$.
\item \emph{(Open bare-LF trace at distinct times)} If $\gamma$ is
self-avoiding and $0 < s_1 < \dots < s_L < \beta$, then
\[
W_{\mathrm{open},\,\mathrm{sync}}(\gamma; s_1, \dots, s_L)
= \exp\!\left[-2 X \coth\!\left(\frac{\beta\Omega}{2}\right)\right]
  \prod_{j=1}^{L-1} F_-(s_{j+1} - s_j)^2.
\]
\item \emph{(Physical closed trace with endpoint observable $\widetilde{\Dop}_{x_0}^\dagger \widetilde{\Dop}_{x_L}$ at $u=0$)}
Including the endpoint displacements $D_{x_0}(-2\alpha)$ and $D_{x_L}(+2\alpha)$
at $u=0$ along a self-avoiding path $\gamma$,
\[
W_{\mathrm{phys},\,\mathrm{sync}}(\gamma; s_1, \dots, s_L)
= F_-(s_1)^2 \left(\prod_{j=1}^{L-1} F_-(s_{j+1} - s_j)^2\right) F_-(s_L)^2.
\]
In particular, as $0 < s_1 < \dots < s_L \to 0^+$, $W_{\mathrm{phys},\,\mathrm{sync}} \to 1$
with no $\exp(-c U/\Omega)$ factor, neither per step nor overall.
\end{enumerate}
\end{theorem}

\begin{proof}
By site factorization and the Gaussian Weyl identity
$\langle \mathcal{T}_u \prod_k D(r_k; u_k)\rangle_{\beta,\mathrm{ph}} = \exp(-\Phi_\beta[q])$,
each intermediate site $x_j$ in~(2) carries charges $-2$ at $s_j$ and $+2$ at $s_{j+1}$,
contributing $\Phi_\beta = \frac{2U}{\Omega}(D_\beta(0) - D_\beta(s_{j+1}-s_j)) = -\log F_-(s_{j+1}-s_j)^2$.
The endpoints contribute
\[
\langle D_{x_0}(+2\alpha)\rangle_{\beta,\mathrm{ph}}
\langle D_{x_L}(-2\alpha)\rangle_{\beta,\mathrm{ph}}
= \exp\!\left[-2X\coth\left(\frac{\beta\Omega}{2}\right)\right]
\]
in~(2) and telescope to $F_-(s_1)^2$ and $F_-(s_L)^2$ against $u=0$ in~(3).
\end{proof}

\begin{theorem}[Near-synchronized time-offset action bound]
\label{thm:tel-offset}
Let an intermediate site $x_j$ be visited once by each spin, entering at
$u_j^\uparrow, u_j^\downarrow \in [0,\beta)$ (charges $-1$ each) and departing
at $u_{j+1}^\uparrow, u_{j+1}^\downarrow \in [0,\beta)$ (charges $+1$ each),
with $\Delta u_j^\sigma := |u_j^\sigma - u_{j+1}^\sigma|_\beta$. Then:
\begin{enumerate}[leftmargin=1.5em]
\item For every $s \in (0,\beta)$,
$|D_\beta'(s)| = \Omega\,|\sinh(\Omega(\beta/2-s))|/\sinh(\beta\Omega/2) \le \Omega$,
so $0 \le D_\beta(0) - D_\beta(s) \le \Omega |s|_\beta$ uniformly in $\beta,\Omega > 0$.
\item The site action satisfies
\begin{equation}
\label{eq:tel-offset}
\Phi_\beta[q_{x_j}] \le U \bigl(\Delta u_j^\uparrow + \Delta u_j^\downarrow\bigr).
\end{equation}
In particular, when $\Delta u_j^\sigma = O(1/U)$, each intermediate step pays
only $O(1)$ phonon penalty rather than $e^{-c U/\Omega}$.
\end{enumerate}
\end{theorem}

\begin{proof}
Writing the charge vector as a sum of two dipoles
$(e_{u_{j+1}^\uparrow} - e_{u_j^\uparrow}) + (e_{u_{j+1}^\downarrow} - e_{u_j^\downarrow})$
in the positive-semidefinite inner product induced by $D_\beta$ and applying
$\|a + b\|_{D_\beta}^2 \le 2(\|a\|_{D_\beta}^2 + \|b\|_{D_\beta}^2)$ together
with $\|e_u - e_v\|_{D_\beta}^2 = 2(D_\beta(0) - D_\beta(|u-v|_\beta)) \le 2\Omega |u-v|_\beta$
from~(1) yields~\eqref{eq:tel-offset}.
\end{proof}

Turning to the fermion determinant in the pre-Lang-Firsov Hubbard-Stratonovich
representation, exact positivity holds without bipartiteness or half-filling.

\begin{lemma}[Generalized Hubbard-Stratonovich positivity and exact dimer cancellation]
\label{lem:tel-hs}
Let $V$ be any finite graph and let $\mu \in \R$ be arbitrary. Integrating out
the local oscillators produces the centered $\beta$-periodic Gaussian process
$\phi = \{\phi_z(u)\}_{z \in V,\, u \in [0,\beta)}$ with covariance
$\E[\phi_z(u)\phi_w(v)] = \delta_{zw} \frac{U\Omega}{2} D_\beta(|u-v|_\beta)$.
For each path $\phi$, let $h_{\phi,\mu}(u) = -t A_V - \diag(\phi_z(u)) - \mu I \in M_{|V|}(\R)$,
$U_{\phi,\mu} = \mathcal{T}_u \exp(-\int_0^\beta h_{\phi,\mu}(u)\,du) = e^{\beta\mu} U_{\phi,0} \in M_{|V|}(\R)$,
$\bar\phi_z = \int_0^\beta \phi_z(u)\,du$, and $G_{\phi,\mu} = (I + U_{\phi,\mu})^{-1} U_{\phi,\mu}$.
Then:
\begin{enumerate}[leftmargin=1.5em]
\item $Z_V(\beta,\mu)/Z_{\mathrm{ph},V}(\beta) = \E_\phi[W_{V,\mu}(\phi)]$
with $W_{V,\mu}(\phi) := (\prod_{z \in V} e^{-\bar\phi_z})\det(I + U_{\phi,\mu})^2 \ge 0$.
\item For any $x, y \in V$,
\[
\langle \Dop_x^\dagger \Dop_y \rangle_{\beta,V,\mu}
= \frac{\E_\phi\bigl[W_{V,\mu}(\phi)\, G_{\phi,\mu}(y,x)^2\bigr]}{\E_\phi[W_{V,\mu}(\phi)]} \ge 0.
\]
\item On the two-site dimer ($V = \{0,1\}$, any $\mu \in \R$), the determinant
cancels identically:
\[
W_{2,\mu}(\phi)\, G_{\phi,\mu}(1,0)^2
= e^{-\bar\phi_0 - \bar\phi_1} U_{\phi,\mu}[1,0]^2
= e^{2\beta\mu - \bar\phi_0 - \bar\phi_1} U_{\phi,0}[1,0]^2 \ge 0.
\]
\end{enumerate}
\end{lemma}

\begin{proof}
Factorizing the two identical spin traces for real $h_{\phi,\mu}(u)$ gives
$\det(I + U_{\phi,\mu})^2 \ge 0$, while $\Dop_x^\dagger \Dop_y = (c_{x\uparrow}^\dagger c_{y\uparrow})(c_{x\downarrow}^\dagger c_{y\downarrow})$
and Wick's theorem yield $\Tr(c_x^\dagger c_y M_{\phi,\mu}^\sigma) = \det(I + U_{\phi,\mu}) G_{\phi,\mu}(y,x) \in \R$.
On $V = \{0,1\}$, the $2\times 2$ adjugate identity $(I + U)^{-1}U = \det(I + U)^{-1}\bigl(\begin{smallmatrix}U_{00}+\det U & U_{01}\\ U_{10} & U_{11}+\det U\end{smallmatrix}\bigr)$
gives $\det(I + U_{\phi,\mu}) G_{\phi,\mu}(1,0) = U_{\phi,\mu}[1,0]$, canceling
$\det(I + U_{\phi,\mu})^2$. (Bipartiteness and $\mu=0$ are needed only for
particle-hole symmetry fixing $\langle n_z\rangle = 1$, not for positivity.)
\end{proof}

Lifting Lemma~\ref{lem:tel-hs} from a finite segment $B_L = \{0,\dots,L\}$ to
a chain $V \supset B_L$ uniformly in $|V|$ gives a boundary self-energy
$\Sigma_{\mathrm{ext}}(u,v)$ supported on $\{0,L\}$ with
$\E_{H_E}[\|\Sigma_{\mathrm{ext}}\|_{\mathrm{HS}}^2] \le \beta^2 t^4$, but a
uniform-in-$|V|$ lower bound
$\E_\phi[W_{V,\mu}(\phi) G_{\phi,\mu,V}(L,0)^2] / \E_\phi[W_{V,\mu}(\phi)] \ge c_d (t/U)^{2d}$
($d=L$) remains open.

\begin{observation}[Exact-diagonalization refutation of a literal per-step physical bound]
\label{obs:tel-ed}
Consider the candidate strong-coupling bound formulated in~\cite{barriers3}
for the physical pair correlator,
\begin{equation}
\label{eq:tel-fg-conj}
\langle \Dop_x^\dagger \Dop_y \rangle_\beta
\le C_0 \left(\frac{C_1 t^2}{U\tau} e^{-c_2 (U/\Omega)\coth(\beta\Omega/4)}\right)^{\lceil d(x,y)/2 \rceil}
  \cdot \mathcal{R}_{\mathrm{neutral}}(\beta, t, U, \Omega),
\end{equation}
with $\mathcal{R}_{\mathrm{neutral}}$ uniformly bounded in $|\Lambda|$. Converged
pre-Lang-Firsov exact diagonalization in the supplement
(\texttt{supp/replay-telescope.py}, all $19$ reported rows having relative
cutoff difference at most $4.92 \times 10^{-3}$) refutes~\eqref{eq:tel-fg-conj}
with $\mathcal{R}_{\mathrm{neutral}} = O(1)$:
\begin{itemize}[leftmargin=1.5em]
\item In the grand-canonical ($\mu=0$) dimer ($N=2, d=1$, converged $M_{\mathrm{tot}} = 40 \to 50$)
at $U=10, \Omega=0.5, t=1.0$ ($X = 20$): for $\tau = 0.25$ ($\beta\Omega = 2$),
$e^{-20\coth(0.5)} = 1.60 \times 10^{-19}$ and
$(t^2/(U\tau)) e^{-20\coth(0.5)} = 6.40 \times 10^{-20}$, whereas exact
diagonalization gives $\langle \Dop_0^\dagger \Dop_1 \rangle_{\mathrm{GC}} = 1.580242 \times 10^{-3} = 0.1580\,(t/U)^2$
(relative cutoff change $2.08 \times 10^{-4}$, $(1,1)$-sector weight $0.6003$),
exceeding the bare kernel factor by more than $10^{16}$. At $\tau = 0.05$
($\beta\Omega = 10$), $e^{-20\coth(2.5)} = 1.57 \times 10^{-9}$, while
$\langle \Dop_0^\dagger \Dop_1 \rangle_{\mathrm{GC}} = 2.294804 \times 10^{-3} = 0.2295\,(t/U)^2$
(relative cutoff change $8.44 \times 10^{-5}$). At $U=12, \Omega=1.0, t=0.25$
($X=12$), $\langle \Dop_0^\dagger \Dop_1 \rangle_{\mathrm{GC}}$ equals
$5.864751 \times 10^{-5} = 0.1351\,(t/U)^2$ at $\tau=0.25$ and
$6.376923 \times 10^{-5} = 0.1469\,(t/U)^2$ at $\tau=1.0$.
\item In canonical $(1,1)$-sector exact diagonalization on $N=2$ ($d=1$) and
$N=3$ ($d=2$), the physical correlator likewise scales as $c_d (t/U)^{2d}$
without per-step $e^{-c U/\Omega}$ suppression (for instance, $0.2598\,(t/U)^2$
on $N=2$ at $U=10,\Omega=0.5,t=1.0,\tau=0.05$, and $7.886803 \times 10^{-8} = 0.0827\,(t/U)^4$
on $N=3$ at $U=8,\Omega=1.0,t=0.25,\tau=0.25$).
\end{itemize}
Consequently, any bound carrying per-step Franck-Condon suppression
$e^{-c (U/\Omega)\coth(\beta\Omega/4)}$ cannot apply directly to the raw physical
correlator $\langle \Dop_x^\dagger \Dop_y \rangle_\beta$ with an $O(1)$
neutral prefactor, and must instead account for virtual short-time dressing or
mechanisms beyond per-step Franck-Condon suppression.
\end{observation}

\section{Annealed spectral estimates and the second-chaos obstruction}
\label{sec:annealed-spectral}

Let $A$ be a real symmetric nearest-neighbor adjacency matrix on a finite graph
of $V$ vertices with absolute row sums at most $q$ ($q = 2d$ on a $d$-dimensional
torus). Set
\[
h_v = -tA + \diag(v) - \mu I, \qquad
C_v = f_\beta(h_v), \qquad
f_\beta(e) = (1 + e^{\beta e})^{-1}, \qquad
\phi_T(e) = T\log(1 + e^{-e/T}),
\]
with $T = 1/\beta$, and define the static free-energy functional
\[
F(v) = \sum_{x=1}^V \frac{v_x^2}{2U} - \sum_{x=1}^V v_x - 2\Tr\phi_T(h_v),
\qquad
\nu_0(dv) = Z_0^{-1} e^{-\beta F(v)}\,dv.
\]
The static annealed probability measure $\nu_0$ is the classical-phonon limit at
fixed $U = 2g^2/\Omega$ as $\Omega \to 0$. In the quantum model, decomposing
the Euclidean field as $v(\tau) = v_0 + u(\tau)$ with
$u(\tau) = \sum_{m \ne 0} u_m e^{i\nu_m \tau}$, $\nu_m = 2\pi m/\beta$,
defines a reference product measure $\nu_0 \times G_>$ where $G_>$ is the
centered Gaussian process with $\E_G[u_{x,m} u_{y,n}] = \delta_{xy}\delta_{m,-n} T \hat{D}(\nu_m)$
and $\hat{D}(\nu) = U\Omega^2/(\nu^2+\Omega^2)$. The full quantum marginal of
$v_0$, obtained by reweighting $\nu_0 \times G_>$ by the squared ratio of
time-dependent to static one-spin determinants, is denoted $\nu_q$.

\begin{theorem}[Static annealed spectral package]
\label{thm:as-static}
For every finite graph $V$ with maximum degree bound $q$ and all $U, T > 0$,
$\mu \in \R$:
\begin{enumerate}[leftmargin=1.5em]
\item \emph{(Static pair sum rule and Kubo bound)} At fixed $v$,
$\langle \Dop_x^\dagger \Dop_y \rangle_v = (C_v)_{xy}^2$,
$\sum_y \langle \Dop_x^\dagger \Dop_y \rangle_v = (C_v^2)_{xx} \le 1$, and
the pair Kubo kernel $K_{xy} = \int_0^\beta \langle \Dop_x^\dagger(\tau)\Dop_y(0)\rangle_v\,d\tau$
satisfies $\sum_y K_{xy} = \sum_a |R_{xa}|^2 \frac{\tanh(\beta e_a/2)}{2e_a} \le \beta/4$,
where $R$ diagonalizes $h_v$ with eigenvalues $\{e_a\}$.
\item \emph{(Uniform-in-$v$ Fermi-kernel decay)} For all $x \ne y$ with $r = d(x,y)$,
\[
|(C_v)_{xy}| \le 4\exp\!\left(-\frac{r}{2(1 + q e \beta|t|/\pi)}\right),
\]
and for $q = 2d$ ($d \ge 1$) one has $|(C_v)_{xy}| \le 4\exp(-a_d r/(1 + \beta|t|))$
with $a_d = \pi/(4de)$.
\item \emph{(Finite-shift relative-entropy level counting)} Under $\nu_0$, for
every $E \ge T$,
\[
N(E) := \frac{1}{V} \E_{\nu_0} \Tr \mathbf{1}_{[-E,E]}(h_v)
\le \min\!\left(1,\, \frac{8E}{U}\right),
\]
with the improved constant $2 + 4\log 2$ in place of $8$ using shift $a = 2(1 + 2\log 2)E$.
\item \emph{(Annealed resolvent and Cooper bubble)} For every $|\omega| \ge T$,
\[
R(\omega) := \frac{1}{V}\E_{\nu_0}\Tr(h_v^2 + \omega^2)^{-1}
\le \min\!\left(\omega^{-2},\, \frac{8(1 + \pi/4)}{U|\omega|}\right),
\]
and the averaged uniform Cooper bubble $B = \frac{1}{V}\Tr[\tanh(\beta h_v/2)/(2h_v)]$
satisfies $\E_{\nu_0} B \le \min\bigl(\frac{1}{4T},\, \frac{4}{U}[1 + \log_+(U/(16T))]\bigr)$.
\item \emph{(Single nonstatic Cooper rung)} With fermionic Matsubara frequencies
$\omega_m = (2m+1)\pi T$ and $K_m(x,y) = |G_v(i\omega_m)_{xy}|^2$,
\[
L_n := \frac{1}{V}\sum_{x,y}\sum_{m \ne n} T\hat{D}(\omega_n - \omega_m) \E_{\nu_0}[K_m(x,y)]
\le \left(\frac{2}{3\pi} + \frac{1}{6}\right)\left(\frac{\Omega}{T}\right)^2,
\]
and the spatially weighted sum with $e^{2b d(x,y)}$ for
$b = \log(1 + \pi T/(2q|t|))$ ($t \ne 0$) gains at most a factor $4$.
\item \emph{(Fisher integration-by-parts identity)} With the two-spin
conditional density susceptibility $\chi_v^\rho(i\nu)$,
\[
\E_{\nu_0} F'' = \beta \E_{\nu_0}[F'(F')^T] \succeq 0,
\qquad
\E_{\nu_0}\chi_v^\rho(0) \preceq U^{-1} I,
\qquad
0 \preceq \chi_v^\rho(i\nu) \preceq \chi_v^\rho(0).
\]
\end{enumerate}
\end{theorem}

\begin{proof}
Part~(1) follows from real-spin Wick factorization and orthonormal diagonalization
of $h_v$. Part~(2) applies the sine representation
$C_{xy} = -\int_0^\infty [\sin(s h_v)]_{xy}/[\beta\sinh(\pi s/\beta)]\,ds$
with interaction-picture Combes-Thomas weighting $W_z = e^{d(z,y)}$; the
$c = qe\beta|t|/\pi \le 1/2$ series bound $\frac{1}{\pi}\int_0^\infty \frac{e^{cu}-1}{\sinh u}\,du \le \pi/4$
and large-$c$ split at $u_* = r/(2(c+1))$ yield prefactor $4$. Part~(3) uses
Lebesgue translation invariance $D(\nu_a \| \nu_0) = \beta \E_{\nu_0}[F(v+a\mathbf{1}) - F(v)] \ge 0$
symmetrized over $\pm a$ ($a = 4E$), canceling linear terms to leave
$2\E_{\nu_0}\Tr[\phi_T(h_v+a) + \phi_T(h_v-a) - 2\phi_T(h_v)] \le V a^2/U$,
where the bracket exceeds $E$ on $[-E,E]$ for $E \ge T$. Layer-cake integration
of $N(E)$ yields~(4), summing $\sum_{m \ne n}(n-m)^{-2}|2m+1|^{-1} \le \pi^2/3$
with $b = \log(1 + \pi T/(2q|t|))$ yields~(5), and Gaussian confinement
justifies componentwise integration by parts for~(6), combined with the
Lehmann representation of $\chi_v^\rho(i\nu)$.
\end{proof}

\begin{theorem}[Continuum first-variation bound]
\label{thm:as-firstvar}
Under $\nu_0 \times G_>$, the linearized equal-time covariance response
$(\delta C)_{ab} = \sum_{m \ne 0} \frac{f_a - f_b}{e_a - e_b + i\nu_m} (u_m)_{ab}$
satisfies
\begin{equation}
\label{eq:as-1over96}
\frac{1}{V} \E_{\nu_0} \E_G \|\delta C\|_{\mathrm{HS}}^2
\le \frac{(\beta\Omega)^2}{96}.
\end{equation}
\end{theorem}

\begin{proof}
For each nonzero mode $\nu_m$, Gaussian averaging $\E_G[u_m u_{-m}] = T\hat{D}(\nu_m)$
together with $|f_a - f_b|^2 \le \frac{\beta}{4}[-(e_a-e_b)(f_a-f_b)]$ and
the factor $2$ for spins in $\chi_v^\rho(i\nu_m)$ gives
$\E_G \|\delta C_m\|_{\mathrm{HS}}^2 \le \frac{\beta}{8} T\hat{D}(\nu_m)\Tr\chi_v^\rho(i\nu_m)$.
Applying Theorem~\ref{thm:as-static}(6) yields
$\frac{1}{V}\E_{\nu_0}\E_G\|\delta C\|_{\mathrm{HS}}^2 \le \frac{1}{8U}\sum_{m\ne 0}\hat{D}(\nu_m)$,
and $\sum_{m\ne 0} \hat{D}(\nu_m) \le U\Omega^2 \sum_{m \ne 0}(2\pi m T)^{-2} = U(\beta\Omega)^2/12$,
giving $\frac{1}{8}\cdot \frac{1}{12} = \frac{1}{96}$.
\end{proof}

\begin{theorem}[Coherent-state pressure sandwich]
\label{thm:as-sandwich}
Define the coherent-state partition function
\[
Z_{\mathrm{coh}}
= \int \prod_{x=1}^V \frac{d^2 z_x}{\pi}\, e^{-\beta\Omega\sum_x |z_x|^2}
  e^{\beta\sum_x v_x} \det\bigl(1 + e^{-\beta h_v}\bigr)^2,
\qquad
v_x = g(z_x + \overline{z_x}).
\]
Then $Z_{\mathrm{coh}} = [\Omega\sqrt{2\pi\beta U}]^{-V} Z_0$ and the full
quantum partition function $Z_q$ satisfies
\begin{equation}
\label{eq:as-sandwich}
Z_{\mathrm{coh}} \le Z_q \le e^{\beta\Omega V} Z_{\mathrm{coh}}.
\end{equation}
\end{theorem}

\begin{proof}
Changing variables to $v_x = 2g\operatorname{Re} z_x$ (Jacobian $1/(2g)$) and
integrating out $\operatorname{Im} z_x$ (factor $\sqrt{\pi/(\beta\Omega)}$)
with $g^2 = U\Omega/2$ yields the prefactor $[\Omega\sqrt{2\pi\beta U}]^{-V}$.
Because lower and upper coherent-state symbols differ by $-\Omega V$, expanding
the coherent-state isometry $T$ in an electronic eigenbasis of the upper
symbol and applying scalar Jensen yields the upper bound, while testing against
$|z\rangle$ tensored with lower-symbol eigenvectors yields the lower bound.
\end{proof}

\begin{theorem}[Quantum zero-mode variance and low-temperature dimer obstruction]
\label{thm:as-dimer}
Under the full quantum marginal $\nu_q$, for every finite graph $V$ and
$U, \Omega, T > 0$,
\begin{equation}
\label{eq:as-nuqvar}
\E_{\nu_q}[v_{0,x}] = -U \langle q_x \rangle_q,
\qquad
\Var_{\nu_q}(v_{0,x}) = UT + U^2 T \chi_{q,xx}^\rho(0),
\end{equation}
where $\chi_{q,xx}^\rho(0) = \int_0^\beta \langle q_x(\tau) q_x(0)\rangle_{q,\mathrm{conn}}\,d\tau$
is the full interacting charge susceptibility. On the half-filled ($\mu=0$)
spinful Holstein dimer ($V=2$) with $t > 0, \Omega > 0, U > 2t$, setting
$U = 16t$ and $E = 3t/2$ gives
\[
\lim_{T \to 0^+} \frac{1}{2} \E_{\nu_q} \Tr \mathbf{1}_{[-E,E]}(h_{v_0}) = 1
> \frac{3}{4} = \frac{8E}{U},
\]
so unrestricted graph-independent transfer of Theorem~\ref{thm:as-static}(3)
from $\nu_0$ to $\nu_q$ fails as $T \to 0$ at fixed $\Omega > 0$ (while $T \ge \Omega$
and fixed $d=2,3$ tori remain open).
\end{theorem}

\begin{proof}
Completing the square with static source $\sum_x J_x v_{0,x}$ in the Euclidean
path integral gives $\E_{\nu_q}[e^{J\cdot v_0}] = e^{\frac{UT}{2}|J|^2}\langle e^{-UT J\cdot Q}\rangle_q$
($Q_x = \int_0^\beta q_x(\tau)\,d\tau$), differentiating twice gives~\eqref{eq:as-nuqvar}.
On the half-filled dimer with $U > 2t$, ground energies in $N=0,4$ equal $-U$
and in $N=1,3$ exceed $-U/2-t > -U$, while in $N_\uparrow = N_\downarrow = 1$
the positivity-improving oscillator heat kernel and connected hopping give a
unique gapped ground state with energy below $-\sqrt{U^2 + t^2} < -U$.
Harmonic confinement gives $\Tr e^{-\beta_0(H-E_0)} < \infty$ at fixed $\beta_0 > 0$,
so excited states contribute $O(\beta e^{-\beta\cdot\mathrm{gap}}) \to 0$ to
$\chi_{q,xx}^\rho(0;T)$, while the ground state contributes at most
$2\|q_x\|^2/\mathrm{gap} < \infty$. Since $\langle q_x \rangle_q = 0$ by
particle-hole symmetry, $\E_{\nu_q}[v_{0,x}^2] = O(T) \to 0$. At $U=16t$,
$h_0$ has eigenvalues $\pm t$, so Weyl's inequality places both eigenvalues of
$h_{v_0}$ inside $[-3t/2, 3t/2]$ whenever $\max_x |v_{0,x}| \le t/2$, and
Markov's inequality drives the level-count expectation to $1 > 3/4$.
\end{proof}

Finally, we analyze the direct second-chaos normalization proposal. On $M$
time slices ($\Delta t = \beta/M$) with $B_k = e^{-\Delta t\,h_v/2} e^{-\Delta t\,\varepsilon \diag(u_k)} e^{-\Delta t\,h_v/2}$
and $P = B_{M-1}\cdots B_0$, write $D = \det(I + P) = D_0(1 + \varepsilon^2 W_2 + O(\varepsilon^3))$
and $C = (I - (I+P)^{-1})^T = C_0 + \varepsilon C_1 + \varepsilon^2 C_2 + O(\varepsilon^3)$.
For a fixed spatial weight $w_{xy} = \exp(a d(x,y))/V$, set $g(v) = \sum_{x,y} w_{xy} (C_0)_{xy}^2$.

\begin{obstruction}[Negative verdict on direct unweighted second-chaos normalization]
\label{obs:as-chaos}
Expanding the normalized annealed pair ratio $\E_{\nu_0 \times G_>}[D^2 \sum_{x,y} w_{xy} C_{xy}^2]/\E_{\nu_0 \times G_>}[D^2]$
in $\varepsilon$ yields:
\begin{enumerate}[leftmargin=1.5em]
\item At order $\varepsilon^2$,
\[
\Delta^{(1)}
= \E_{\nu_0}\sum_{x,y} w_{xy}\bigl[\E_G C_{1,xy}^2 + 2 C_{0,xy} \E_G C_{2,xy}\bigr]
+ 2\Cov_{\nu_0}(g, \E_G W_2),
\]
with continuum mean $\E_G W_2 = \frac{1}{4}\sum_{\nu \ne 0}\hat{D}(\nu)\Tr\chi_v^\rho(i\nu)$,
while at leading pure second-Hermite-chaos order $\varepsilon^4$,
\[
\Delta_{H_2}
= \E_{\nu_0}\sum_{x,y} w_{xy}\bigl[\Var_G C_{2,xy} + 2 C_{0,xy}\Cov_G(C_{2,xy}, W_2)\bigr]
+ \Cov_{\nu_0}(g, \Var_G W_2).
\]
Neither normalization term is an unweighted expectation of $-\E[|\nabla F|^2]$,
and $W_2$ probes nonzero-frequency $\chi_v^\rho(i\nu)$ rather than $\chi_v^\rho(0)$.
\item Under the static weighted integration-by-parts identity
\begin{equation}
\label{eq:as-ward}
\Cov_{\nu_0}\bigl(g, \Tr\chi^\rho(0)\bigr)
= -\beta\Cov_{\nu_0}\bigl(g, |\nabla F|^2\bigr) + \E_{\nu_0}\bigl[\nabla g \cdot \nabla F\bigr],
\end{equation}
the gradient term $\E_{\nu_0}[\nabla g \cdot \nabla F]$ cannot be dropped:
on the one-site $t=0, \mu=0$ model with $g(v) = f_\beta(v)^2$, as $U \to 0$
one has $\nu_0 \to \delta_0$ and $\E_{\nu_0}[g' F'] = \E_{\nu_0}[g'']/\beta \to \beta/8 \ne 0$
(and by continuity for small $t > 0$ on connected graphs). Furthermore, in
\texttt{supp/replay-annealed-spectral.py} ($t=1, U=2, \beta=3, \Omega=0.1, a=1$),
the normalization term $2\Cov_{\nu_0}(g, \E_G W_2)/(\beta\Omega)^2$ is strictly
positive across all three tested graphs ($+3.132\times 10^{-4}$ on the dimer,
$+2.900\times 10^{-4}$ on the four-site ring, and $+2.587\times 10^{-5}$ on the
$2\times 2$ periodic torus), and in the static Ward check ($t=1, U=2, \beta=2.5, \mu=0.3, a=1$)
both $-\beta\Cov(g, |\nabla F|^2) = +4.670\times 10^{-3}$ and $\E[\nabla g \cdot \nabla F] = -1.797\times 10^{-2}$
are essential.
\end{enumerate}
Hence direct combination of the normalization covariance into the positive
unweighted static force-square of Theorem~\ref{thm:as-static}(6) fails; any
successful all-orders Gaussian-chaos reorganization remains open.
\end{obstruction}

\section{Annealed endpoint reduction, static Combes-Thomas base, and instanton annihilation}
\label{sec:annealed-endpoint}

We now turn to the annealed endpoint formulation~\cite{fisher,capstone} on a
connected finite simple graph $\Lambda$ of $V = |\Lambda|$ vertices and
maximum degree $d_{\max} \ge 1$. Write $\tau = \beta^{-1}$,
$\epsilon = (\Omega/\tau)^2 = (\beta\Omega)^2$, $\lambda = U\epsilon/\tau$,
and decompose the Matsubara field as $b(u) = b_0 + \sqrt{s}\, b_{\mathrm{dyn}}(u)$
with $s \in [0,1]$, where $b_{\mathrm{dyn}}$ has centered Gaussian Matsubara
modes $\E_{\mathrm{dyn}}[|b_{n,z}|^2] = U\tau \epsilon/(4\pi^2 n^2 + \epsilon)$
($n \ne 0$). Set $h(u) = -tA - \mu I_V - \diag(b_0 + \sqrt{s}\,b_{\mathrm{dyn}}(u))$
(note the opposite sign convention on $b_0$ relative to $v_0$ in
Section~\ref{sec:annealed-spectral}), monodromy
$U_{\mathrm{tot}}[b] = \mathcal{T}\exp(-\int_0^\beta h(u)\,du)$,
occupied Green's matrix $G_{\mathrm{occ}} = U_{\mathrm{tot}}(I_V + U_{\mathrm{tot}})^{-1}$,
and $Q = (I_V + U_{\mathrm{tot}})^{-1}$. Conditional on $b_0$ and $s$, denote
$S_s(b_0) = \frac{1}{V}\langle \Tr(G_{\mathrm{occ}}^T G_{\mathrm{occ}}) \rangle_{b_0,s}$
and $\tilde{S}_s(b_0) = \frac{1}{V}\langle \Tr(Q^T Q) \rangle_{b_0,s}$, and
write $\nu_1(db_0) = p_{\mathrm{dyn}}(b_0)\,db_0 \propto e^{-\beta |b_0|^2/(2U)} Z_{\mathrm{cond}}(b_0)\,db_0$
for the exact $s=1$ static-mode marginal, with conditional charge susceptibility
$(\chi_{b_0,1}^{\mathrm{cond}})_{xy} = \tau\frac{\partial^2 \log Z_{\mathrm{cond}}(b_0)}{\partial b_{0,x}\partial b_{0,y}}$.

In the notation of~\cite{fisher,capstone}, the annealed endpoint conjecture
posits that for $\epsilon \le \epsilon_0(t,U)$,
\begin{equation}
\label{eq:ae-conj1}
\E_{\nu_1}\bigl[S_1(b_0) - 1\bigr]
\le A(t,U)\,\epsilon\,\frac{U}{V}\,\E_{\nu_1}\bigl[\Tr \chi_{b_0,1}^{\mathrm{cond}}\bigr],
\end{equation}
whereas the spatially weighted closing estimate posits that for $t > 0$ and
Combes-Thomas mass $m(\tau, t, d_{\max}) = 2\log(1 + \frac{\pi\tau}{2 d_{\max} t}) > 0$,
the weighted moment $\mathcal{P}_m := \sup_{x\in\Lambda}\sum_{y\in\Lambda} e^{m d(x,y)} \langle \Delta_x^\dagger \Delta_y\rangle$
satisfies, for all $\epsilon \le \epsilon_0(t,U)$,
\begin{equation}
\label{eq:ae-closing2}
\mathcal{P}_m \le K_0 + A(t,U)\,\epsilon\,\mathcal{P}_m
\end{equation}
with constants $K_0, A(t,U)$ independent of $V$ and $\mu$.

\begin{theorem}[Conditional reduction and insufficiency of unweighted summability]
\label{thm:ae-cond}
Let $\Lambda$ be any connected finite simple graph with maximum degree at most
$d_{\max} \ge 1$.
\begin{enumerate}[leftmargin=1.5em]
\item Combining~\eqref{eq:ae-conj1} with the exact interacting conditional
Fisher identity~\cite{fisher}
$\E_{\nu_1}[\chi_{b_0,1}^{\mathrm{cond}}] + \tau^{-1}\E_{\nu_1}[(\nabla_{b_0}\Phi_{\mathrm{dyn},1})(\nabla_{b_0}\Phi_{\mathrm{dyn},1})^T] = U^{-1}I_V$
yields $\frac{U}{V}\E_{\nu_1}[\Tr \chi_{b_0,1}^{\mathrm{cond}}] \le 1$ and hence
$\E_{\nu_1}[S_1] \le 1 + A(t,U)\epsilon$ uniformly in $V$ and $\mu$.
\item Inequality~\eqref{eq:ae-conj1} alone yields only $\ell^1$ pair
summability ($\frac{1}{V}\sum_{x,y}\langle\Delta_x^\dagger \Delta_y\rangle \le 1 + A\epsilon$)
and does not imply exponential decay~\eqref{eq:tpair-def}: on a $d$-dimensional
torus, a translation-invariant power-law profile $c_0(1+d(x,y))^{-(d+\eta)}$
($\eta > 0$) has bounded sum uniformly in $L$ yet violates $C q^{d(x,y)}$ for
every $q \in (0,1)$.
\item For $t > 0$, if Closing Estimate~\eqref{eq:ae-closing2} holds at mass
$m(\tau,t,d_{\max})$ for all $\epsilon \le \epsilon_0(t,U)$ with $V$- and
$\mu$-uniform constants $K_0, A(t,U) < \infty$, then setting
$\epsilon_*(t,U) := \min\{\epsilon_0(t,U), (2A(t,U))^{-1}\}$ and
$C(t,U) := \epsilon_*(t,U)^{-1/2} = \max\{\epsilon_0(t,U)^{-1/2}, \sqrt{2A(t,U)}\}$
forces $\mathcal{P}_m \le 2 K_0$ whenever $\tau > C(t,U)\Omega$, so
$0 \le \langle \Delta_x^\dagger \Delta_y \rangle \le 2 K_0 q^{d(x,y)}$ with
$q = (1 + \frac{\pi\tau}{2 d_{\max} t})^{-2} < 1$, and therefore
$k_B \Tpair \le C(t,U)\Omega$. (At $t = 0$, off-diagonal correlators vanish
identically so $\Tpair = 0$.)
\end{enumerate}
\end{theorem}

\begin{proof}
Part~(1) takes the trace of the conditional Fisher identity and uses
positive semidefiniteness of $(\nabla_{b_0}\Phi_{\mathrm{dyn},1})(\nabla_{b_0}\Phi_{\mathrm{dyn},1})^T \succeq 0$.
Part~(2) follows directly from $\sum_{r=1}^\infty r^{d-1}(1+r)^{-(d+\eta)} < \infty$
versus $\lim_{r\to\infty}(1+r)^{-(d+\eta)} q^{-r} = +\infty$. For~(3), if
$\tau > C(t,U)\Omega$ then $\epsilon < \epsilon_*(t,U) \le (2A(t,U))^{-1}$;
because $\mathcal{P}_m \le V e^{mV} < \infty$ on any finite connected graph
and every pair correlator is nonnegative, rearranging~\eqref{eq:ae-closing2}
gives $\mathcal{P}_m \le K_0/(1 - A(t,U)\epsilon) \le 2K_0$. Markov's inequality
then yields $e^{m d(x,y)}\langle \Delta_x^\dagger \Delta_y\rangle \le 2K_0$.
\end{proof}

\begin{theorem}[Universal static-limit Combes-Thomas base bound]
\label{thm:ae-k0}
At $s = 0$ on any connected finite simple graph $\Lambda$ of maximum degree
$d_{\max} \ge 1$ and for all $\tau > 0, \mu \in \R, b_0 \in \R^V$:
\begin{enumerate}[leftmargin=1.5em]
\item At $t = 0$, off-diagonal pair correlators vanish identically and
$\mathcal{P}_m^{(0)}(b_0)|_{t=0} = \max_x f_\tau(-\mu-b_{0,x})^2 < 1$ for
every finite $m \ge 0$.
\item For every $t > 0$, with $m(\tau,t,d_{\max}) = 2\log(1 + \frac{\pi\tau}{2 d_{\max} t})$,
both static pair correlators satisfy
\begin{equation}
\label{eq:ae-k0}
\mathcal{P}_m^{(0)}(b_0)
\le K_0 := \left(1 + \frac{\pi}{4}\right)^2 \in [3.1876466018,\, 3.1876466019]
< 1 + \frac{\pi^2}{4},
\end{equation}
uniformly in $b_0, \mu, V, t > 0$, and $\tau > 0$.
\end{enumerate}
In addition, \texttt{supp/replay-annealed-endpoint.py} verifies the outward
interval Schur row and column test on a six-site ring at equality rows and
columns ($\|B_\eta\|_{\ell^1 \to \ell^1} \le \pi\tau/2$, with normalized row~$3$
and column~$0$ sums enclosed in $[2,2]$) at $60$-digit \texttt{mpmath.iv} precision.
\end{theorem}

\begin{proof}
Fix $x_0 \in \Lambda$ and set $\eta = m/2 = \log(1 + \frac{\pi\tau}{2 d_{\max} t})$,
so $e^\eta - 1 = \pi\tau/(2 d_{\max} t)$. For $(D_\eta \psi)(y) = e^{\eta d(x_0,y)}\psi(y)$~\cite{combes1973},
symmetry of $h_{b_0}$ gives
$\sum_y e^{m d(x_0,y)} (f_\tau(h_{b_0})_{yx_0})^2 \le \|D_\eta f_\tau(h_{b_0}) D_\eta^{-1}\|_{\mathrm{op}}^2$.
Using the Matsubara partial-fraction expansion
$f_\tau(h_{b_0}) = \frac{1}{2}I_V - \tau \sum_{k\in\Z}(i\omega_k I_V + h_{b_0})^{-1}$
($\omega_k = (2k+1)\pi\tau$), conjugation leaves the diagonal potential
unchanged and shifts $h_{b_0}$ by $(B_\eta)_{yz} = -t A_{yz}(e^{\eta(d(x_0,y)-d(x_0,z))} - 1)$.
Since $|d(x_0,y)-d(x_0,z)| \le 1$ on edges and row/column degrees are at most
$d_{\max}$, Schur's test bounds $\|B_\eta\|_{\mathrm{op}} \le d_{\max} t(e^\eta - 1) = \pi\tau/2$.
Self-adjointness of $h_{b_0}$ gives $\|R_k\|_{\mathrm{op}} = \|(i\omega_k I_V + h_{b_0})^{-1}\|_{\mathrm{op}} \le |\omega_k|^{-1}$,
so $\|R_k B_\eta\|_{\mathrm{op}} \le \frac{1}{2|2k+1|} \le \frac{1}{2}$ and
$\|D_\eta R_k D_\eta^{-1} - R_k\|_{\mathrm{op}} \le 2\|B_\eta\|_{\mathrm{op}}\|R_k\|_{\mathrm{op}}^2 \le (\pi\tau (2k+1)^2)^{-1}$.
Summing over $k \in \Z$ gives $\|D_\eta f_\tau(h_{b_0}) D_\eta^{-1} - f_\tau(h_{b_0})\|_{\mathrm{op}} \le \frac{2}{\pi}\sum_{k=0}^\infty (2k+1)^{-2} = \pi/4$,
so $\|D_\eta f_\tau(h_{b_0}) D_\eta^{-1}\|_{\mathrm{op}} \le 1 + \pi/4$, and
similarly for $Q_{b_0} = I_V - f_\tau(h_{b_0})$.
\end{proof}

Before turning to the $t=0$ instanton analysis and unconditional $V=2$ proof,
we record that $S_p[b] < 1$ and $\tilde{S}_p[b] < 1$ hold pointwise whenever
$U_{\mathrm{tot}}[b]$ is normal with $\operatorname{Re}\lambda_j > 0$ for every
eigenvalue (covering static $s=0$, atomic $t=0$, and time-reversal-even paths
$b_{\mathrm{dyn}}(\beta - u) = b_{\mathrm{dyn}}(u)$ where $U_{\mathrm{tot}} = U_{1/2}^T U_{1/2} \succ 0$),
whereas a $2\times 2$ orthogonal rotation $R(3\pi/4)$ has $\det(I_2 + R) = 2-\sqrt{2} > 0$
yet $S_p = 1 + 1/\sqrt{2} > 1$; time-translation invariance of $Z_{\mathrm{el}}[b]$
averages $S_p$ over time shifts, and Gaussian integration by parts governs
dynamic modes with variance $U\tau\epsilon/(4\pi^2 n^2 + \epsilon)$.

\begin{theorem}[Instanton suppression and unconditional $V=2$ proof]
\label{thm:ae-v2}
Consider the two-site dimer ($V=2$) with $t \ge 0$, $U, \Omega, \tau > 0$.
\begin{enumerate}[leftmargin=1.5em]
\item \emph{(Exact $t=0$ nine-component marginal at $\mu=0$)} At $t=0$ and $\mu=0$,
$Z_{\mathrm{cond}}(b_0)|_{t=0} = \prod_{x=1}^2 (e^{\beta b_{0,x}} + 2 + e^{-\beta b_{0,x}}) = 16\cosh^2(\beta b_{0,1}/2)\cosh^2(\beta b_{0,2}/2)$,
so multiplying by $e^{-\beta(b_{0,1}^2+b_{0,2}^2)/(2U)}$ and completing the square
makes $p_{\mathrm{dyn}}(b_0)|_{t=0}$ a product of two three-component Gaussian
mixtures of variance $U\tau$ centered at $\{-U, 0, +U\}$ (bare coefficients $1:2:1$,
completed-square weights $e^{\beta U/2}:2:e^{\beta U/2}$), i.e., a nine-component
mixture on $\{-U, 0, +U\}^2$. (In rotated coordinates $b_{0,\pm} = (b_{0,1}\pm b_{0,2})/\sqrt{2}$,
$Z_{\mathrm{cond}}|_{t=0} = 4(\cosh(\beta b_{0,+}/\sqrt{2}) + \cosh(\beta b_{0,-}/\sqrt{2}))^2$
does not factorize.)
\item \emph{(Exact $t=0$ barrier-to-well ratio)} For every $B \ge 0$ and $\beta, U > 0$,
the barrier-to-well pointwise marginal density ratio at $b_0^{\mathrm{bar}} = (-B, +B)$ satisfies
\begin{equation}
\label{eq:ae-r0}
\begin{aligned}
R_0(B; \beta, U)
&:= \frac{p_{\mathrm{dyn}}((-B,+B))|_{t=0}}{p_{\mathrm{dyn}}((-U,+U))|_{t=0}}
= \left(\frac{1 + e^{-\beta B}}{1 + e^{-\beta U}}\right)^4 \exp\!\left(-\frac{\beta(U-B)^2}{U}\right) \\
&\le C_{\mathrm{bar}}(\beta U)\,\exp\!\left(-\frac{\beta(U-B)^2}{U}\right),
\end{aligned}
\end{equation}
where $C_{\mathrm{bar}}(\beta U) := 16(1 + e^{-\beta U})^{-4} \in (1, 16)$.
Its exact supremum over $B \in [0, U]$ (and $[0,\infty)$) is attained at
$B_*(\beta U) = 0$ when $\beta U \le 2$ (value $16 e^{-\beta U}(1 + e^{-\beta U})^{-4}$)
and at the unique positive root $B_*(\beta U) \in (0, U)$ of
$\tanh(\beta B_*/2) = B_*/U$ when $\beta U > 2$. In particular, for separated
barriers $U - B \ge \sqrt{U\tau\log(16/\delta)}$, $R_0 \le \delta \ll 1$
(at the instanton point $U=3.0, B=1.0, \tau=0.1$, $R_0 \approx 1.620\times 10^{-6}$),
while in the unseparated regime $\beta U = 1, B = 0.9U$ one can have $R_0 \approx 1.1069 > 1$.
\item \emph{(Unconditional $V=2$ operator-norm proof)} For every $t \ge 0$,
$U, \Omega, \tau > 0$, and $\mu \in \R$ at $V = 2$,
\[
\E_{\nu_1}[S_1(b_0)] \le 1 - P(0) - P(1) - \frac{1}{2}P(3) < 1,
\qquad
\E_{\nu_1}[\tilde{S}_1(b_0)] \le 1 - \frac{1}{2}P(1) - P(3) - P(4) < 1,
\]
where $P(N) = \langle \mathbf{1}_{\mathcal{F}_N}\rangle > 0$ is the exact thermal
weight of the $N$-electron sector $\mathcal{F}_N$. Thus~\eqref{eq:ae-conj1} holds
unconditionally at $V=2$ with $A = 0$ for all $\epsilon > 0$.
\end{enumerate}
\end{theorem}

\begin{proof}
Parts~(1) and~(2) follow from $U_{\mathrm{tot}}|_{t=0} = \diag(e^{\beta b_{0,1}}, e^{\beta b_{0,2}})$,
writing
\[
w_0(\pm B) = e^{\beta U/2} e^{-\beta(U-B)^2/(2U)}(1 + e^{-\beta B})^2,
\]
and observing that $\frac{\partial}{\partial B}\log w_0(B) = \beta(\tanh(\beta B/2) - B/U)$
is strictly concave on $[0, \infty)$ with derivative $\frac{\beta}{2} - \frac{1}{U}$ at $0$.
For~(3), on the $16$-dimensional dimer Fock space $\bigoplus_{N=0}^4 \mathcal{F}_N$,
$\hat{S}_1 = \frac{1}{2}(\Delta_1^\dagger + \Delta_2^\dagger)(\Delta_1 + \Delta_2)$
commutes with $N$ and phonon operators and has sector operator norms
$(0, 0, 1, 1/2, 1)$ on $\mathcal{F}_0,\dots,\mathcal{F}_4$ (and $(1, 1/2, 1, 0, 0)$
for $\hat{\tilde{S}}_1$), as verified in exact rational arithmetic by the
supplement script. Since $P(N) > 0$ and $\sum_{N=0}^4 P(N) = 1$,
taking the thermal expectation yields the two strict bounds. (For $V \ge 4$,
by contrast, the Anderson pseudospin $SU(2)$ maximum eigenvalue of $\hat{S}_1$
is $V^{-1}\lfloor (V+1)/2\rfloor \lceil (V+1)/2\rceil \ge 3/2 > 1$.)
\end{proof}

\begin{theorem}[Static strictness and exact first-order $s=0$ derivative]
\label{thm:ae-s0}
Fix a connected finite simple graph $\Lambda$ ($V = |\Lambda|$), $t, U, \tau > 0$,
and $\mu \in \R$.
\begin{enumerate}[leftmargin=1.5em]
\item At $s = 0$, $1 - S_0(b_0) = \frac{1}{V}\sum_{a=1}^V q_a(1+f_a) \ge \frac{1}{V}\sum_{a=1}^V f_a q_a > 0$
and $\Tr\chi_{b_0,0}^{\mathrm{cond}} \ge \frac{2\beta}{V}\sum_{a=1}^V f_a q_a > 0$
($f_a = f_\tau(\varepsilon_a(b_0))$, $q_a = 1 - f_a$). At $\mu = 0$ on any
bipartite graph, $\E_{\nu_0}[S_0(b_0)] = \frac{1}{2} - \frac{1}{V}\E_{\nu_0}[\sum_{a=1}^V f_a q_a] < \frac{1}{2}$.
\item At fixed $(\tau, t, U, \mu, V)$, $\E_{\nu_1}[S_1 - 1] = -\E_{\nu_0}[1 - S_0] + O(\epsilon) < 0$
and $\frac{U}{V}\E_{\nu_1}[\Tr\chi_{b_0,1}^{\mathrm{cond}}] > 0$ for all $\epsilon \in [0, \epsilon_1(\tau,t,U,\mu,V)]$.
\item By~\cite{fisher} Theorem~C.1, the exact first-order dynamic derivative
at $s = 0$ satisfies, for every $b_0 \in \R^V$,
\begin{equation}
\label{eq:ae-c1phys}
\begin{aligned}
\left|\left.\frac{\partial \tilde{S}_s(b_0)}{\partial s}\right|_{s=0}\right|
&\le C_1^{\mathrm{phys}}\,\frac{U\Omega^2}{\tau^2}\left(\frac{1}{V}\sum_{z\in\Lambda} (\chi_{b_0,0})_{zz}\right)
\le \frac{C_1^{\mathrm{phys}}}{2}\,\lambda, \\
C_1^{\mathrm{phys}} &:= \frac{1}{24} + \frac{\zeta(3)}{8\pi^3} \in [0.0465126891,\, 0.0465126892],
\end{aligned}
\end{equation}
with $\frac{1}{2}C_1^{\mathrm{phys}} \in [0.0232563445,\, 0.0232563446]$ and
$\E_{\nu_0}[|\partial_s \tilde{S}_s|_{s=0}|] \le C_1^{\mathrm{phys}}\epsilon$.
\end{enumerate}
\end{theorem}

\begin{proof}
Part~(1) follows from $f_a + q_a = 1$, Cauchy-Schwarz $V\Tr\chi_{b_0,0}^{\mathrm{cond}} \ge \mathbf{1}^T \chi_{b_0,0}^{\mathrm{cond}}\mathbf{1}$,
and on bipartite graphs $\Sigma h_{b_0}\Sigma = -h_{-b_0}$ at $\mu=0$, which swaps
$S_0(-b_0) = \tilde{S}_0(b_0)$ under the even measure $\nu_0$. Part~(2) follows
from smoothness of the Gaussian covariance in $\epsilon$, and Part~(3) combines
\cite{fisher} Theorem~C.1 with $0 \le (\chi_{b_0,0})_{zz} \le (2\tau)^{-1}$ and
$\frac{U}{V}\E_{\nu_0}[\Tr\chi_{b_0,0}] \le 1$.
\end{proof}

\begin{proposition}[High-frequency attractive Hubbard anchor and global barriers]
\label{prop:ae-hubbard}
On any fixed finite graph $\Lambda$ at fixed $(t, U, \tau, \mu)$, as $\Omega \to \infty$
the Holstein partition function and equal-time correlators converge to those of
the attractive Hubbard Hamiltonian $H_{\mathrm{Hub}} = -t\sum_{\langle x,y\rangle,\sigma}(c_{x\sigma}^\dagger c_{y\sigma} + \mathrm{h.c.}) - \mu\sum_x n_x - \frac{U}{2}\sum_x(n_x - 1)^2$.
On bipartite graphs at $t=1, U=4, \tau=0.1, \mu=0$, Hubbard exact diagonalization
gives $\langle S_1\rangle_{\mathrm{Hub}} = 0.853160 < 1$ ($V=2$), $1.048208 > 1$ ($V=4$),
and $1.058314 > 1$ ($V=6$), increasing further for $\mu > 0$ (e.g., $1.318580$ at
$V=6, \mu=0.5$). Meanwhile, numerical Fock-ED and exact-covariance UDT-DQMC checks
at $V = 2, 4, 6$ confirm $\E_{\nu_1}[S_1] < 1$ across all tested $\epsilon \le 64$
regimes (including the pathwise barrier points where $S_1(b_0^{\mathrm{bar}}) = 20.51$ at $V=2$,
$3.86$ at $V=4$, and $2.40$ at $V=6$). Two explicit mathematical barriers keep a
uniform finite-$s$ proof of~\eqref{eq:ae-conj1} and~\eqref{eq:ae-closing2} open
as $V \to \infty, \tau \to 0$:
\begin{itemize}[leftmargin=1.5em]
\item \emph{(Barrier B1: extensive dynamic covariance)} Differentiating
$\langle S_s(b_0)\rangle$ in $s \in (0,1)$ generates $\Cov_{b_0,s}(S_p, -\partial_s H_{\mathrm{dyn}})$,
which scales as $\Theta(V\epsilon)$ in phase-coexistence static backgrounds $b_0$.
\item \emph{(Barrier B2: finite-$\Omega$ off-diagonal pair coherence)} Separating
the retarded $\Omega \ll U$ regime (where pair hopping suffers Franck-Condon
suppression $\sim (2t^2/U)e^{-U/\Omega}$) from the $\Omega \gtrsim U$ Hubbard
regime (where $\langle S_1\rangle > 1$ on $V \ge 4$) uniformly as $\tau \to 0$
and $V \to \infty$ requires a non-perturbative multi-site bound not yet established.
\end{itemize}
\end{proposition}

\section{Volume-uniform neutral-dipole cluster expansion and collar boundary}
\label{sec:neutral-dipole}
 
We finally isolate the neutral-sector ultraviolet/infrared resummation and its
interface with a charged skeleton $\mathcal{C}$. Let $\Lambda$ be a finite simple
graph of maximum vertex degree $\Delta \ge 1$ ($q = \Delta$; $\Delta = 2d$ on a
$d$-dimensional torus). Fix a static occupation configuration $\eta \in \{0,2\}^\Lambda$
(two electrons or zero electrons per site), and define the neutral-loop partition
function $Z_{\mathrm{neut}}^{\mathrm{loop}}(\Lambda; \eta) = Z_{\mathrm{el},\eta}^{-1} \Tr_{\mathrm{neut}}(\eta)\exp(-\beta H_{\mathrm{LF}})$
on the continuous-time spacetime $\Lambda \times S^1_\beta$. A neutral polymer
$\Gamma$ is a connected monochromatic cycle of single-spin hops $c_{z_r\sigma}^\dagger c_{z_{r-1}\sigma}$
returning every site $z \in V(\Gamma)$ to $\eta_z$ without changing the local
pair state on any interval; its active set $A_z(\Gamma) = \{s \in S^1_\beta : n_z^{\Gamma}(s) = 1\}$
is a disjoint union of open intervals with total active length
$L(\Gamma) = \sum_{z \in V(\Gamma)} |A_z(\Gamma)|$. Two polymers $\Gamma_1, \Gamma_2$
are incompatible ($\Gamma_1 \sim \Gamma_2$) if $A_z(\Gamma_1) \cap A_z(\Gamma_2) \ne \emptyset$
for some $z \in \Lambda$, and compatible ($\Gamma_1 \not\sim \Gamma_2$) otherwise.

For a compatible collection $\{\Gamma_i\}$, the exact phonon action splits as
$\Phi_\beta(\bigcup_i \Gamma_i) = \sum_i \Phi_\beta(\Gamma_i) + \sum_{i < j} V(\Gamma_i, \Gamma_j)$,
where $V(\Gamma, \Gamma') = \sum_{z \in V(\Gamma) \cap V(\Gamma')} V_z(A_z(\Gamma), A_z(\Gamma'))$
is the dipole-dipole phonon cross-interaction
\[
V_z(I, J) = -\frac{U\Omega}{2} \int_I \int_J \delta q_z(I)\delta q_z(J)\, D_\beta(s - s')\,ds\,ds',
\qquad
\delta q_z(\Gamma) = 1 - \eta_z \in \{-1, +1\}.
\]
Because $\Phi_\beta \ge 0$ unconditionally, compatibility satisfies the
Poghosyan-Ueltschi~\cite{poghosyan2009} stability condition (Assumption~1)
without truncation, while the integrated two-body kernel satisfies the exact
time-translation bound $\int_{\text{compatible } \tau} |V_z(\Gamma, \Gamma' + \tau)|\,d\tau \le U |A_z(\Gamma)| |A_z(\Gamma')|$
uniformly in $\beta, \Omega, \mu$, and $\eta$. Writing the signed polymer measure as
\[
d\mu_\eta(\Gamma) = \operatorname{sgn}_\eta(\Gamma)\,
\frac{t^{m(\Gamma)}}{m(\Gamma)}\,
e^{-\frac{U}{2}L(\Gamma)-\Phi_\beta(\Gamma)}\,
\mathbf{1}_{\{\Gamma\ \eta\text{-legal}\}}\,ds_1\cdots ds_{m(\Gamma)}.
\]
Here $1/m(\Gamma)$ accounts for cyclic relabeling of the ordered hop times;
the factor $e^{-\frac{U}{2}L(\Gamma)}$ comes from the atomic excitation penalty
$+\frac{U}{2}$ for $n_z = 1$ relative to $n_z \in \{0,2\}$, while $V_z \le 0$ and
convexity of $D_\beta$ on $[0,\beta/2]$ yield the uniform phonon-action bounds
$0 \le \Phi_\beta(\Gamma) \le \frac{U}{2} L(\Gamma)$ (with $\Phi_{\beta,z}([0,u]) \to \frac{U}{2\Omega}$ for $u, \beta - u \gg \Omega^{-1}$).

\begin{theorem}[Volume-, $\beta$-, $\Omega$-, $\mu$-, and $\eta$-uniform cluster convergence]
\label{thm:nd-conv}
Define the Poghosyan-Ueltschi test functions~\cite{poghosyan2009} (extending Koteck\'y-Preiss~\cite{kotecky1986}) $a(\Gamma) = a_0 m(\Gamma) + a_1 U L(\Gamma)$
and $b(\Gamma) = \Phi_\beta(\Gamma) \ge 0$, with auxiliary parameter $\varepsilon > 0$
and $c := \frac{1}{2} - a_1 - \varepsilon > 0$.
Whenever $t/U \le c_0(\Delta)$ (with explicit graph-family thresholds in
Table~\ref{tab:nd-constants}), the neutral-loop polymer measure $d\mu_\eta$
satisfies the Poghosyan-Ueltschi~\cite{poghosyan2009} Assumption~2' pinning
criterion
\[
\int_{\mathcal{P}_{\mathrm{neut}}(\eta)} d|\mu_\eta|(\Gamma')\, e^{a(\Gamma') + b(\Gamma')}
\bigl(1_{\{\Gamma' \sim \Gamma\}} + 1_{\{\Gamma' \not\sim \Gamma\}} |V(\Gamma, \Gamma')|\bigr)
\le a(\Gamma) \qquad \forall \Gamma \in \mathcal{P}_{\mathrm{neut}}(\eta),
\]
uniformly in $|\Lambda|, \beta > 0, \Omega > 0, \mu \in \R$, and $\eta \in \{0,2\}^\Lambda$.
The separate qualitative hypothesis $\int d|\mu_\eta|(\Gamma)\,
e^{a(\Gamma)+2b(\Gamma)} < \infty$ holds for each finite $\Lambda,\beta,\Omega$;
no uniform bound on that integral is needed.
\end{theorem}

\begin{proof}
Because $\Phi_\beta(\bigcup_i \Gamma_i) \ge 0$, stability (Assumption~1) holds
with $b(\Gamma) = \Phi_\beta(\Gamma)$, while finite $\Lambda, \beta, \Omega$ give
$\int d|\mu_\eta|(\Gamma)\,e^{a(\Gamma)+2b(\Gamma)} < \infty$. In Assumption~2',
the choice $b(\Gamma') = \Phi_\beta(\Gamma')$ cancels $e^{-\Phi_\beta(\Gamma')}$
identically, leaving $e^{-(\frac{1}{2}-a_1) U L(\Gamma') + a_0 m(\Gamma')}$ in
$d|\mu_\eta|(\Gamma')\,e^{a(\Gamma')+b(\Gamma')}$. Rooting each $m$-step polymer
at a chosen $(z_0, s_1)$ fixes the spin $\sigma$ and gives at most
$2(2\Delta)^{m-1}$ rooted shapes (and $4\Delta$ for $m=2$; on bipartite graphs
odd $m$ are absent).
Integrating over the global time shift first bounds the hard-core overlap
measure by $\frac{m(\Gamma')}{2}|A_z(\Gamma)| + k_z(\Gamma)|A_z(\Gamma')|$,
where $k_z$ counts active intervals. The compatible soft term obeys
\[
\int_{\text{compatible }\tau}|V(\Gamma,\Gamma'+\tau)|\,d\tau
\le U\sum_z |A_z(\Gamma)|\,|A_z(\Gamma')|.
\]
Using $x e^{-\varepsilon x} \le (e\varepsilon)^{-1}$ for
$x=U L(\Gamma')$ changes the remaining decay rate to
$c=\frac{1}{2}-a_1-\varepsilon$. Integration over the $m-1$ active interval
lengths then gives $S_L U L(\Gamma)+S_m\sum_z k_z(\Gamma)$,
which is at most $a(\Gamma)$ since $\sum_z k_z(\Gamma)=m(\Gamma)$,
$S_L\le a_1$, and $S_m\le a_0$ for $t/U\le c_0(\Delta)$.
\end{proof}

\begin{corollary}[Pure-neutral pressure analyticity and site pinning]
\label{cor:nd-press}
For $t/U \le c_0(\Delta)$, the finite-volume neutral pressure
$p_{\mathrm{neut}}(\Lambda; \eta) := \frac{1}{\beta |\Lambda|}\log Z_{\mathrm{neut}}^{\mathrm{loop}}(\Lambda; \eta)$
converges absolutely and uniformly in $|\Lambda|, \beta, \Omega, \mu, \eta$, with
\[
\bigl|p_{\mathrm{neut}}(\Lambda; \eta)\bigr|
\le \frac{1}{\beta |\Lambda|} \int_{\mathcal{P}_{\mathrm{neut}}(\eta)} d|\mu_\eta|(\Gamma)\, e^{a(\Gamma) + b(\Gamma)}
\le C_{\mathrm{press}}(\Delta)\,\frac{t^2}{U},
\]
and for any measurable subset $\mathcal{R} \subset \mathcal{P}_{\mathrm{neut}}(\eta)$,
the cluster sum over clusters touching $\mathcal{R}$ is bounded by
$\int_{\mathcal{R}} d|\mu_\eta|(\Gamma_1)\, e^{a(\Gamma_1) + b(\Gamma_1)}$.
\end{corollary}

\begin{lemma}[Unperturbed neutral-sector ratio bounds]
\label{lem:nd-ratios}
For $t/U \le c_0(\Delta)$:
\begin{enumerate}[leftmargin=1.5em]
\item \emph{(Background difference)} For any $\eta, \eta' \in \{0,2\}^\Lambda$
differing on $D = \{z \in \Lambda : \eta_z \ne \eta'_z\}$,
\[
\left|\log \frac{Z_{\mathrm{neut}}^{\mathrm{loop}}(\Lambda; \eta)}{Z_{\mathrm{neut}}^{\mathrm{loop}}(\Lambda; \eta')}\right|
\le C_{\mathrm{diff}}(\Delta)\, |D|\,\beta\,\frac{t^2}{U},
\qquad
C_{\mathrm{diff}}(\Delta) = 2 C_{\mathrm{press}}(\Delta).
\]
\item \emph{(Spacetime exclusion)} For any $S = \bigsqcup_{z\in\Lambda}(\{z\} \times S_z)$
with $S_z \subset S^1_\beta$ consisting of $k_z(S)$ connected intervals,
\[
\left|\log \frac{Z_{\mathrm{neut}}^{\mathrm{loop}}(\Lambda; \eta)}{Z_{\mathrm{neut}}^{\mathrm{loop}}(\Lambda \setminus S; \eta)}\right|
\le \sum_{z\in\Lambda} \left(C_{\mathrm{press}}(\Delta)\, |S_z|\,\frac{t^2}{U} + C_{\mathrm{col},2}(\Delta)\, k_z(S)\,\frac{t^2}{U^2}\right).
\]
\end{enumerate}
\end{lemma}

\begin{proof}
In~(1), all clusters disjoint from $D$ cancel identically between $\eta$ and
$\eta'$, and applying Corollary~\ref{cor:nd-press} to clusters touching $D$
gives $2 C_{\mathrm{press}}(\Delta)|D|\beta t^2/U$. In~(2), clusters disjoint
from $S$ cancel; for a rooted polymer touching $S_z$, the translation measure
of $s_1$ giving $I_1 \cap S_z \ne \emptyset$ is at most $|S_z| + k_z(S)L(\Gamma'_0)$,
and bounding $L e^{-c^* U L} \le \frac{1}{e\varepsilon' U} e^{-(c^*-\varepsilon')UL}$
($c^* = \frac{1}{2} - a_1$, $\varepsilon' = 0.10$) yields $C_{\mathrm{col},2}(\Delta)$.
\end{proof}

To couple neutral loops to an arbitrary charged skeleton $\mathcal{C}$ whose
charged segments occupy $R_{\mathcal{C}} = \bigsqcup_z (\{z\}\times R_{\mathcal{C},z})$
($k_{\mathrm{seg},z}(\mathcal{C})$ closed intervals) with locally constant
occupation $\eta_{\mathrm{loc}}(z, s) \in \{0,2\}$ outside $R_{\mathcal{C}}$,
define the $\Omega^{-1}$-collar $R_{\mathrm{col}} = \bigsqcup_z (\{z\}\times R_{\mathrm{col},z})$
by $R_{\mathrm{col},z} = \{s \in S^1_\beta : \operatorname{dist}_{S^1_\beta}(s, R_{\mathcal{C},z}) \le w_{\mathrm{col}}\}$
with $w_{\mathrm{col}} := (\ln 40)/\Omega$ (so $\delta_{\mathrm{col}} := 2e^{-\Omega w_{\mathrm{col}}} = 0.05$)
and $k_{\mathrm{col},z} \le k_{\mathrm{seg},z}(\mathcal{C})$ connected intervals.
On the exterior $\Lambda \times S^1_\beta \setminus R_{\mathrm{col}}$, let
$Z_{\mathrm{neut}}^{\mathrm{ext}}(\mathcal{C}; R_{\mathrm{col}})$ denote the
cluster expansion over neutral cycles $\Gamma \in \mathcal{P}_{\mathrm{ext}}(\mathcal{C}; R_{\mathrm{col}})$
with weight $e^{-V(\mathcal{C},\Gamma)}\,d\mu_{\eta_{\mathrm{loc}}}(\Gamma)$,
where $V(\mathcal{C},\Gamma) = \sum_{z, I \subset A_z(\Gamma)} V_z(\mathcal{C}, I)$
is the exact phonon cross-interaction with $\mathcal{C}$.

\begin{theorem}[Perturbed exterior cluster expansion and collar-ratio interface]
\label{thm:nd-ext}
Let $\mathcal{C}$ be any charged-skeleton history on $\Lambda \times S^1_\beta$,
and fix any reference background $\eta \in \{0,2\}^\Lambda$. Let
$D_{\mathcal{C},z}^{\mathrm{st}} = \{s \in S^1_\beta \setminus R_{\mathrm{col},z} : \eta_{\mathrm{loc}}(z,s) \ne \eta_z\}$
have $k_{\mathrm{flip},z} \le k_{\mathrm{seg},z}(\mathcal{C})$ connected intervals.
\begin{enumerate}[leftmargin=1.5em]
\item \emph{(Exterior fermion sign factorization)} In the full trace with
$\mathcal{C}$ and exterior neutral loops $\Gamma_1,\dots,\Gamma_n \in \mathcal{P}_{\mathrm{ext}}(\mathcal{C}; R_{\mathrm{col}})$,
$\sgn(\mathcal{C} \cup \Gamma_1 \cup \dots \cup \Gamma_n) = \sgn(\mathcal{C})\prod_{i=1}^n \sgn_{\eta_{\mathrm{loc}}}(\Gamma_i)$,
where $\sgn_{\eta_{\mathrm{loc}}}(\Gamma) = (-1)^{N_2^{\mathrm{loc}}(\Gamma)-1}$.
\item \emph{(Pointwise exterior potential bound)} For every $z \in \Lambda$ and
$s \in S^1_\beta \setminus R_{\mathcal{C},z}$, the phonon cross-field obeys
$\bigl|\frac{U\Omega}{2}\int_{S^1_\beta}(n_z^{\mathcal{C}}(s') - n_z^{\mathcal{C}}(s)) D_\beta(s-s')\,ds'\bigr| \le 2U e^{-\Omega\operatorname{dist}_{S^1_\beta}(s, R_{\mathcal{C},z})}$.
Hence $|V(\mathcal{C},\Gamma)| \le \delta_{\mathrm{col}} U L(\Gamma) = 0.05\, U L(\Gamma)$
for all $\Gamma \in \mathcal{P}_{\mathrm{ext}}(\mathcal{C}; R_{\mathrm{col}})$, and
the exterior $L^1$ integral is at most $0.10\,(U/\Omega)\,k_{\mathrm{seg},z}(\mathcal{C})$.
\item \emph{(Perturbed PU convergence)} For $t/U \le c_0^{\mathrm{ext}}(\Delta)$
(replacing $\frac{1}{2}-a_1-\varepsilon$ by $c_{\mathrm{ext}} = \frac{1}{2}-\delta_{\mathrm{col}}-a_1-\varepsilon$),
the interpolated exterior cluster expansion at $\lambda \in [0,1]$ converges
uniformly in $\mathcal{C}, |\Lambda|, \beta, \Omega, \mu$, and $\lambda$.
\item \emph{(Exact collar-ratio interface)} For $t/U \le c_0^{\mathrm{ext}}(\Delta)$,
uniformly in $\Lambda, \beta, \Omega, \mu, \eta$, and $\mathcal{C}$,
\begin{multline*}
\left|\log \frac{Z_{\mathrm{neut}}^{\mathrm{ext}}(\mathcal{C}; R_{\mathrm{col}})}{Z_{\mathrm{neut}}^{\mathrm{loop}}(\Lambda; \eta)}\right|
\le \sum_{z \in \Lambda} \Bigl(
C_{\mathrm{ext}}(\Delta)\, k_{\mathrm{seg},z}(\mathcal{C})\,\frac{t^2}{U\Omega} \\
+ C_{\mathrm{diff}}(\Delta)\, |D_{\mathcal{C},z}^{\mathrm{st}}|\,\frac{t^2}{U}
+ C_{\mathrm{press}}(\Delta)\, |R_{\mathrm{col},z}|\,\frac{t^2}{U}
+ C_{\mathrm{col},2}(\Delta)\, (2 k_{\mathrm{flip},z} + k_{\mathrm{col},z})\,\frac{t^2}{U^2}
\Bigr),
\end{multline*}
with explicit constants listed in Table~\ref{tab:nd-constants}.
\end{enumerate}
\end{theorem}

\begin{proof}
Fermion-even contraction on each connected component of $S^1_\beta \setminus R_{\mathrm{col},z}$
proves~(1). Because $n_z^{\mathcal{C}}(s') = n_z^{\mathcal{C}}(s)$ for
$|s'-s|_\beta < d = \operatorname{dist}_{S^1_\beta}(s, R_{\mathcal{C},z})$ and
$|n_z^{\mathcal{C}}(s') - n_z^{\mathcal{C}}(s)| \le 2$, integrating $2U\Omega D_\beta(r)$
from $d$ to $\beta/2$ yields $2U\sinh(\Omega(\beta/2-d))/\sinh(\beta\Omega/2) \le 2U e^{-\Omega d}$,
proving~(2). In~(3), Lemma~\ref{lem:nd-ratios} and Theorem~\ref{thm:nd-conv}'s
stability and double-FD bounds remain intact even when intervals in different
components of $S^1_\beta \setminus R_{\mathrm{col},z}$ carry opposite signs,
while $e^{-\lambda V(\mathcal{C},\Gamma)} \le e^{\delta_{\mathrm{col}} U L(\Gamma)}$
shifts $\frac{1}{2} - a_1$ to $c_*^{\mathrm{ext}} = \frac{1}{2} - \delta_{\mathrm{col}} - a_1$.
In~(4), decomposing the log ratio into $(\lambda=1 \text{ vs } \lambda=0)$ in
the exterior, $(\eta_{\mathrm{loc}} \text{ vs } \eta)$ on $S^1_\beta \setminus R_{\mathrm{col}}$,
and $(\Lambda \setminus R_{\mathrm{col}} \text{ vs } \Lambda)$ combines Lemma~\ref{lem:nd-ratios}
with $\lambda$-differentiation and Tonelli exchange against the $O(U/\Omega)$
exterior integral of~(2).
\end{proof}

\begin{table}[t]
\centering
\caption{Uniform cluster-expansion thresholds and constants from Theorems~\ref{thm:nd-conv} and~\ref{thm:nd-ext}, Corollary~\ref{cor:nd-press}, and Lemma~\ref{lem:nd-ratios}.}
\label{tab:nd-constants}
\small
\begin{tabular}{@{}lcccccc@{}}
\toprule
Graph family ($\Delta$) & $c_0(\Delta)$ & $C_{\mathrm{press}}(\Delta)$ & $C_{\mathrm{diff}}(\Delta)$ & $C_{\mathrm{col},2}(\Delta)$ & $c_0^{\mathrm{ext}}(\Delta)$ & $C_{\mathrm{ext}}(\Delta)$ \\
\midrule
$\Delta = 2$ bipartite ($1$D chain) & $0.0294$ & $8.1317$ & $16.2633$ & $49.7231$ & $0.0257$ & $2.8738$ \\
$\Delta = 4$ bipartite ($2$D square) & $0.0181$ & $18.4003$ & $36.8006$ & $126.4678$ & $0.0159$ & $6.9129$ \\
$\Delta = 4$ general & $0.0151$ & $26.7837$ & $53.5674$ & $181.1853$ & $0.0133$ & $7.2104$ \\
$\Delta = 6$ general ($3$D cubic) & $0.0113$ & $44.4016$ & $88.8032$ & $318.0380$ & $0.0099$ & $12.2577$ \\
\bottomrule
\end{tabular}
\end{table}

\begin{remark}[Honest collar boundary and virtual pair-return channels]
\label{rem:nd-scope}
Theorem~\ref{thm:nd-ext} eliminates the extensive $\beta t^2/U$ boundary mismatch
of Obstruction~\ref{obs:cs-boundary} on the exterior $\Lambda \times S^1_\beta \setminus R_{\mathrm{col}}$
whenever $\eta_{\mathrm{loc}} = \eta$ (and bounds mismatches on $D_{\mathcal{C}}^{\mathrm{st}}$
at $C_{\mathrm{diff}} |D_{\mathcal{C}}^{\mathrm{st}}| t^2/U$), leaving only
$O(t^2/(U\Omega) + t^2/U^2)$ collar terms per segment endpoint. Two distinct
mechanisms remain explicitly outside the scope of Theorem~\ref{thm:nd-ext}:
(i)~loops $\Gamma$ whose active set intersects the $\Omega^{-1}$-collar
$R_{\mathrm{col}}$ ($A(\Gamma) \cap R_{\mathrm{col}} \ne \emptyset$), where
near a $\Delta q = \pm 2$ charged jump an attractive screening shift
$V(\mathcal{C},\Gamma) < 0$ produces an $e^{+c U/\Omega}$ boundary-screening
enhancement (reaching $-0.3667\, U/\Omega$ for a two-site dipole at $\beta\Omega=2, X=40$
and $-0.386\, U/\Omega$ to $-0.448\, U/\Omega$ in optimized collar geometries
in \texttt{supp/replay-neutral-dipole.py}); and (ii)~four-hop virtual
pair-return dipoles $(2,0) \to (1,1) \to (0,2) \to (1,1) \to (2,0)$ visiting
the opposite paired state for duration $\Delta s \sim \Omega^{-1}$, which
generate an $O(\beta t^4/(U^2\Omega))$ energy shift without $e^{-U/\Omega}$
suppression. Because persistent background shifts on $D_{\mathcal{C}}^{\mathrm{st}}$
incur $C_{\mathrm{diff}} |D_{\mathcal{C}}^{\mathrm{st}}| t^2/U$, any full-Gibbs
application of Theorem~\ref{thm:nd-ext} operates in the window $U/\Omega \gg \beta t^2/U$
($\tau \gg (t/U)^2 \Omega$) unless paired channels provide further cancellation.
\end{remark}
 
\section{Scope of verification, open problems, and disclosures}
\label{sec:scope}
 
 \paragraph{What is proved, what is ruled out, and what remains open.}
 Every theorem, lemma, proposition, and corollary in this paper is accompanied
 by a proof, and every route obstruction is supported by an explicit mathematical
counterexample or exact identity. Table~\ref{tab:scope} summarizes the
status of each program component:
\begin{itemize}[leftmargin=1.5em]
\item \emph{Section~\ref{sec:charged-skeleton}:} Establishes the Parseval
identity, disjoint-bump Gram inequality, and additive multi-transition plateau
barrier ($3/22$ and $3/19$) in full Fock space, while ruling out naive signless
$|\sigma|$-majorization ($16/160$ negative four-hop neutral loops) and
unexpanded-bath boundary subtraction at $\beta t^2/U \gg 1$.
\item \emph{Section~\ref{sec:telescope}:} Proves the synchronized-pair
telescoping obstruction, generalized Hubbard-Stratonovich determinant
positivity for all $\mu \in \R$ with exact $N=2$ dimer cancellation, and
exact-diagonalization refutation of raw correlators carrying per-step
$e^{-cU/\Omega}$ decay with $O(1)$ neutral prefactor; a volume-uniform Step~2
lower bound remains open.
\item \emph{Section~\ref{sec:annealed-spectral}:} Establishes the static
annealed spectral package, the $(\beta\Omega)^2/96$ continuum first-variation
bound, the coherent-state pressure sandwich, and the quantum zero-mode variance
identity, while ruling out unrestricted graph-independent $\nu_0 \to \nu_q$
level-count transfer as $T \to 0$ (dimer $U=16t, E=3t/2$) and ruling out
direct unweighted second-chaos normalization via the weighted Ward identity.
\item \emph{Section~\ref{sec:annealed-endpoint}:} Proves the conditional
closing reduction (Theorem~\ref{thm:ae-cond}), the universal static base
bound $K_0 = (1+\pi/4)^2$ (Theorem~\ref{thm:ae-k0}), the $t=0$ nine-component
marginal and barrier-to-well suppression $R_0 \le C_{\mathrm{bar}}(\beta U) e^{-\beta(U-B)^2/U}$
plus unconditional $V=2$ operator-norm proof (Theorem~\ref{thm:ae-v2}), and
the exact $s=0$ derivative constant $C_1^{\mathrm{phys}} = 1/24 + \zeta(3)/(8\pi^3)$
(Theorem~\ref{thm:ae-s0}), while leaving uniform finite-$s$ closure and
$\Omega \ll U, \tau \to 0$ multi-site pair coherence open (Barriers~B1 and~B2).
\item \emph{Section~\ref{sec:neutral-dipole}:} Proves the volume-, $\beta$-,
$\Omega$-, $\mu$-, and $\eta$-uniform neutral-loop cluster expansion for
$t/U \le c_0(\Delta)$ and the perturbed exterior resummation for
$t/U \le c_0^{\mathrm{ext}}(\Delta)$, while explicitly leaving open the
$e^{+c U/\Omega}$ attractive collar screening inside $R_{\mathrm{col}}$ and
four-hop virtual pair-return channels of order $\beta t^4/(U^2\Omega)$.
\end{itemize}
The target global phonon-frequency ceiling $k_B \Tpair \le C(t,U)\Omega$
(possibly up to logarithmic factors) in the full unprojected spinful Holstein
Gibbs state remains open.

\begin{table}[t]
\centering
\caption{Summary of analytical results, route obstructions, and open gaps across Sections~\ref{sec:charged-skeleton}--\ref{sec:neutral-dipole}.}
\label{tab:scope}
\small
\begin{tabular}{@{}>{\raggedright\arraybackslash}p{2.9cm}>{\raggedright\arraybackslash}p{4.3cm}>{\raggedright\arraybackslash}p{4.1cm}>{\raggedright\arraybackslash}p{3.2cm}@{}}
\toprule
Component & Established rigorously & Ruled out / negative verdict & Remaining open gap \\
\midrule
Charged skeleton (Sec.~\ref{sec:charged-skeleton}) & Parseval identity, Gram $G \le D$, $3/22$ and $3/19$ multi-transition barriers & Signless $|\sigma|$-majorization ($16/160$ loops); unexpanded-bath subtraction at $\tau \sim \Omega \ll t^2/U$ & Interface decoration outside $S$ across $\eta_{\Lambda\setminus S}$ \\
Telescoping \& ED (Sec.~\ref{sec:telescope}) & Endpoint telescoping cap; $\det^2 \ge 0$ for all $\mu\in\R$ and $N=2$ cancellation & Per-step $e^{-cX}$ bound on raw correlators with $O(1)$ neutral prefactor & Volume-uniform Step~2 determinant ratio lower bound \\
Annealed spectral (Sec.~\ref{sec:annealed-spectral}) & Static $\nu_0$ package; $(\beta\Omega)^2/96$ bound; $Z_{\mathrm{coh}}$ sandwich; $\Var_{\nu_q}(v_0)$ & Graph-independent $\nu_0 \to \nu_q$ level counting as $T\to 0$; direct unweighted $H_2$ force-square identity & Nonlinear weighted $H_2$ chaos bound; $\nu_q$ transfer at $T \ge \Omega$ \\
Annealed endpoint (Sec.~\ref{sec:annealed-endpoint}) & Closing reduction~\eqref{eq:ae-closing2}; $K_0 = (1+\pi/4)^2$; $R_0$ bound; $V=2$ $\E[S_1]<1$; $C_1^{\mathrm{phys}}$ & Pathwise endpoint inequality at staggered barriers; unweighted $S_1 \le 1$ sufficiency for $\Tpair$ & Barriers~B1--B2: uniform finite-$s$ closure and multi-site $\Omega \ll U$ coherence \\
Neutral dipole (Sec.~\ref{sec:neutral-dipole}) & Uniform PU convergence ($t/U \le c_0(\Delta)$); $p_{\mathrm{neut}}$, ratios, exterior collar resummation & Persistence mismatch away from $R_{\mathrm{col}}$ when $\eta_{\mathrm{loc}} = \eta$ & Collar $e^{+c U/\Omega}$ screening inside $R_{\mathrm{col}}$; $O(\beta t^4/(U^2\Omega))$ pair returns \\
\bottomrule
\end{tabular}
\end{table}
\paragraph{Replay supplement.}
Five standalone Python scripts in the supplement (\texttt{supp/}) reproduce all
finite-arithmetic, exact-rational, outward-interval, exact-diagonalization,
determinant quantum Monte Carlo, and worldline Monte Carlo checks referenced in
this paper, and verify $45$ deliberately broken mutant variants across the five
components ($8+7+7+13+10$). Each script exits with status zero when run
standalone from \texttt{supp/}, and \texttt{supp/run\_all.py} runs the complete
battery. Passing these finite checks verifies finite arithmetic and finite-size
witnesses only; it proves no thermodynamic-limit theorem.

\paragraph{Relation to released results.}
Nothing in this paper invalidates or alters any released theorem in the
program \cite{barriers1,barriers2,barriers3,largeu,strong,linear5,fisher,capstone}.
Where Section~\ref{sec:telescope} refutes a literal strong-coupling ceiling, it
refutes the unproved conjecture recorded in Section~3 of the preliminary
companion summary underlying \cite{barriers3}, not any theorem claimed as proved
there.

\paragraph{Use of software and language models.}
AI tools were used for drafting, calculation and checking under the author's
direction and review. All theorems were checked by hand against the proofs, and
all numerical claims are reproducible from the scripts described above.

This preprint has not been peer reviewed.

\paragraph{Data availability.}
The manuscript source, the bibliography and the five replay scripts are archived
with this preprint. Released program results are cited by DOI
\cite{barriers1,barriers2,barriers3,largeu,strong,linear5,fisher,capstone,optimum,finiteenergy,stage6}.

\bibliographystyle{plain}
\bibliography{refs}

\end{document}
